% !TeX program = lualatex
\documentclass[11pt,a4paper]{ltjsarticle}
\usepackage[margin=25mm,headheight=17pt]{geometry}
\usepackage[haranoaji]{luatexja-preset}
\usepackage{amsmath,amssymb,amsthm,mathtools,mathrsfs,booktabs,array,enumitem,longtable}
\usepackage[unicode,hidelinks]{hyperref}
\usepackage{fancyhdr}
\pagestyle{fancy}\fancyhf{}\fancyfoot[C]{\thepage}
\fancyhead[L]{\small\nouppercase{\leftmark}}
\renewcommand{\sectionmark}[1]{\markboth{#1}{}}
\setlength{\emergencystretch}{3em}\allowdisplaybreaks[2]
\setlength{\parskip}{.15em}
\newtheoremstyle{jp}{.6em}{.6em}{\normalfont}{}{\bfseries}{.}{.5em}{}
\theoremstyle{jp}
\newtheorem{theorem}{定理}[section]
\newtheorem{proposition}[theorem]{命題}
\newtheorem{lemma}[theorem]{補題}
\newtheorem{corollary}[theorem]{系}
\newtheorem{definition}[theorem]{定義}
\newtheorem{remark}[theorem]{注意}
\newtheorem{example}[theorem]{例}
\newtheorem{question}[theorem]{問題}
\renewcommand{\proofname}{証明}
\newcommand{\kk}{\mathbb C}\newcommand{\PP}{\mathbb P}
\newcommand{\OO}{\mathcal O}\newcommand{\A}{\mathcal A}
\newcommand{\B}{B_{\mathrm{VK}}}\newcommand{\Blo}{B_{p,q}}
\newcommand{\HH}{\operatorname{HH}}\newcommand{\Hom}{\operatorname{Hom}}
\newcommand{\End}{\operatorname{End}}\newcommand{\Sym}{\operatorname{Sym}}
\newcommand{\Perf}{\operatorname{Perf}}\newcommand{\Der}{\operatorname{Der}}
\newcommand{\Inn}{\operatorname{Inn}}\newcommand{\gr}{\operatorname{gr}}
\newcommand{\rk}{\operatorname{rank}}\newcommand{\Spec}{\operatorname{Spec}}
\newcommand{\id}{\operatorname{id}}\newcommand{\Ext}{\operatorname{Ext}}
\newcommand{\Tor}{\operatorname{Tor}}\newcommand{\Ann}{\operatorname{Ann}}
\newcommand{\coker}{\operatorname{coker}}\newcommand{\Supp}{\operatorname{Supp}}
\newcommand{\pd}{\operatorname{pd}}\newcommand{\depth}{\operatorname{depth}}
\newcommand{\Frac}{\operatorname{Frac}}\newcommand{\Aut}{\operatorname{Aut}}
\newcommand{\Out}{\operatorname{Out}}\newcommand{\GL}{\operatorname{GL}}
\newcommand{\PGL}{\operatorname{PGL}}\newcommand{\Hilb}{\operatorname{Hilb}}
\newcommand{\dashsim}{\mathrel{\overset{\mathrm{bir}}{\sim}}}
\setcounter{tocdepth}{1}

\hypersetup{pdftitle={種数三Prymモジュライの相対有理性と二次被覆},pdfauthor={}}
\title{種数三Prymモジュライの相対有理性と\\二次被覆\\[.4em]
\large 関係ファイバーの指数六と明示的なAlbanese写像}
\author{}\date{2026年10月29日\quad 第30版}
\begin{document}
\maketitle
\begin{abstract}
種数三Prymモジュライから二次曲線のnetを経て得る古典的写像
$\Psi:\mathcal R_3\dashrightarrow Y_0(2)$ を、一般netの有限安定化群を残して調べる。
位数十八の安定化群の基礎体上の形にも適用できる不変式計算とベクトル束の降下により、
$\Psi$ の一般ファイバーがモジュラー関数体上で有理であることを示す。
従ってこの写像を保った双有理同型
$\mathcal R_3\dashsim Y_0(2)\times\mathbb P^5$ が存在する。
一方、Prymモジュライ全体の関数体上では、Hermite曲線の周期と指数は三、
Prym線束の降下障害を表す二次曲線の指数は二である。
自己反対Fano72型関係代数の一般ファイバーは両者の積と双有理であり、
その指数と良い開部分への最小多重切断次数は六となる。
さらに、六次元交代形式の外積三次縮約から階数二のtwisted空間を作り、
そのPl\"ucker二次曲線上の次数二因子によって内在的な二次分裂体を得る。
この体は二つの共通等方三平面を区別する体であり、
関係ファイバーに作用する一次元トーラスの分裂体でもある。
さらに、逆Gram構成が一般の交代二形式の三つ組を合同変換を許して支配することを示す。
これにより、固定net上でも行列式線束がJacobian内で非定数に動く。
共通平面の交換とnet安定化群の符号を掛け合わせた第三の二次被覆について、
双対二次曲線の係数の三次pencilからKummer写像を明示する。
幾何学的一般net上で、この被覆のAlbanese曲線はnetの判別三次曲線のJacobianであり、
不正則数はちょうど一である。Albanese写像の幾何学的一般ファイバーは四次元で単有理である。
行列式線束が与えるHermiteのJacobianへの写像は、その後に二次同種写像を合成したものとなる。
従って後者の幾何学的一般ファイバーは二成分を持つ。
\end{abstract}
\begingroup\small\setlength{\parskip}{0pt}\tableofcontents\endgroup

\section{序論と主結果}

モジュライ空間の有理性と、そこで自然に定まる曲線や線束の降下は、異なる情報を持つ。
種数三Prymモジュライ $\mathcal R_3$ の有理性はDolgachev\cite{DolgachevRationality2008}によって知られている。
同論文は、Prym対の対称二次式表示から二次曲線のnetを取り、
非零二捩れを持つ楕円曲線のモジュライ $Y_0(2)$ へ写す構成も与えている。
ただし一般netには位数十八の安定化群が残るので、
netを固定した二次曲線の五次元空間をそのまま粗モジュライのファイバーとみなすことはできない。

本稿は、この安定化群を保った商の関数体を計算する。
その結果、既知の有理性を、指定されたnet写像に関する相対的な純超越性へ精密化する。
同時に、別稿\cite{PaperV}で関係代数のファイバーと結び付けたHermite曲線に対し、
粗モジュライの関数体上での周期、指数、Brauer障害を決定する。
さらに、その障害を消す二次拡大を、共通等方三平面のPl\"ucker表示から内在的に取り出す。
三つの二次被覆のうち、二つの符号の積で定まる被覆には不正則数一の構造が現れる。
本稿は、そのAlbanese曲線とHermiteのJacobianの間の二次同種写像を明示する。

以下で $C/K$ を一般Prym曲線、$\eta\in\operatorname{Pic}^0_{C/K}(K)[2]\setminus\{0\}$ を
二捩れclassとし、$\alpha=\delta_C(\eta)$ を線束の降下障害とする。
net $N$ のHermite曲線を $E$、$\alpha$ の二次曲線を $D_\alpha$ と書く。
それらの定義、および関係代数の一般ファイバー $S$ への対応は
第\ref{sec:vi-setup-v28}節で述べる。

\begin{theorem}[相対有理性と一般ファイバーの算術]
\label{thm:vi-main-v28}
一般点について次が成立する。
\begin{enumerate}
\item $F=\mathbb C(Y_0(2))$ とすると、古典的net写像 $\Psi$ の一般ファイバーの関数体は
$F$ の五変数純超越拡大である。従って $Y_0(2)$ 上で
\[
 \mathcal R_3\dashsim Y_0(2)\times\mathbb P^5.
\]
\item $K=\mathbb C(\mathcal R_3)$ 上で
\[
 \operatorname{per}(E)=\operatorname{ind}(E)=3,\qquad
 \operatorname{per}(\alpha)=\operatorname{ind}(\alpha)=2.
\]
自己反対関係族の一般ファイバーの滑らかな射影モデルは
\[
 S\dashsim_K E\times_KD_\alpha,\qquad \operatorname{ind}(S)=6
\]
を満たす。その任意の非空開部分に次数六の閉点が存在し、
正則な一般ファイバーを取る良い有理商の最小多重切断次数は六である。
\item 階数六のtwisted空間 $T$ 上の三つの交代形式の縮約
\[
 c_\Omega:\bigwedge^3T\longrightarrow T\otimes\bigwedge^2V
\]
の核 $\Pi$ は階数二であり、$\mathbb P(\Pi)$ はBrauer class $\alpha$ の二次曲線へ降りる。
そのPl\"ucker二次関係式が定める次数二因子は
$\operatorname{Spec}L$ という非自明な二次拡大 $L/K$ である。
これは共通等方三平面の二点を区別する体で、$\alpha_L=0$ を満たす。
\item $L$ 上では $\mathcal K$ が二つの次数 $-3$ の線束に分かれ、
\[
 S_L\dashsim_L E_L\times\mathbb P^1_L,\qquad
 \operatorname{per}(E_L)=\operatorname{ind}(E_L)=\operatorname{ind}(S_L)=3.
\]
また $\operatorname{Res}_{L/K}\mathbb G_m/\mathbb G_m$ が $S\to E$ に忠実に作用し、
固定部分は $E_L\to E$ という次数二の有限エタール多重切断となる。
\item 逆Gram族の合同飽和は、六次元空間上の交代二形式の三つ組の四十五次元空間を支配する。
従って $(E,\{\theta,\theta^{-1}\})$ の二次元幾何学的モジュライを支配し、
一般netを固定しても $\theta$ のKummer座標は非定数である。
\item 一般netの安定化群の符号に付随する二次拡大を $K_\chi=K(\sqrt\epsilon)$、
共通平面の拡大を $L=K(\sqrt\delta)$ と書く。
この二指標は独立で、合成体は次数四である。
第三中間体 $K(\sqrt{\delta\epsilon})$ の模型 $Z$ は、固定したnetのJacobianへ支配的に写る。
$Y_0(2)$ 上の幾何学的一般ファイバーの任意の滑らかな射影モデルは $q=1$ であり、
単有理でなく、安定有理でもない。
\item 幾何学的一般net $N$ 上で、第三被覆 $Z$ のAlbaneseは
$B\simeq\operatorname{Jac}(\Delta_N)$ である。Albanese写像 $Z\dashrightarrow B$ の
幾何学的一般ファイバーは四次元で、幾何学的に整かつ単有理である。
HermiteのJacobian $J$ への写像は $Z\dashrightarrow B\xrightarrow{\varphi}J$ と分解し、
$\varphi$ は次数二の同種写像である。従って $Z\dashrightarrow J$ の幾何学的一般ファイバーは二成分となる。
有限なHesse型パラメータ $m$ 上で、双対二次曲線係数の三次式 $I_m,D$ による
$R=2I_m/D$ を用いると
\[
 B:\ v^2=-R\bigl(4R^2-(m^2-40m-32)R-8(m-2)^3\bigr),
\]
\[
 x_J=-4R+m^2-40m-32+\frac{8(m-2)^3}{R}
\]
と書ける。元の共通平面被覆 $M_L$ の幾何学的一般netファイバーは $q=0$ である。
\end{enumerate}
\end{theorem}

第一項は第\ref{sec:modular-fiber-rationality-v28}節で証明する。
一般netの安定化群の $\operatorname{SL}_3$ への逆像を符号指標でねじると、
位数五十四の単項行列群になる。その不変式環は次数三、三、六の多項式環である。
二次式表現が三次元表現の二重化となることを用い、
射影商を重み付き射影平面上の階数三ベクトル束へ帰着する。
この議論では、安定化群が基礎体上の定数群であることを仮定しない。

第二項の指数三は、支配的な三次捩れ族の離散付値から得る。
指数二は、非分裂二次曲線から作るPrym四次曲線の支配族によって検出する。
両者が互いに素であることが、関係ファイバーの指数六を与える。
第三・四項では、二十次元から十八次元への線形縮約と
六次元空間の同時無限小symplectic自己準同型を用い、
二次分裂代数とトーラスを同じ内在的データから構成する。

第一項は、$\Psi$ が全定義域上で射影空間束となることを主張しない。
第二項の多重切断は、滑らかな射影モデルまたは正則な良い開部分に関する。
任意の特異な境界へ向かう切断を禁止するという意味ではない。
さらに、補助楕円曲線の存在から、一般の二次元Prym多様体に楕円因子があるとは結論しない。

\section{Prym対、netと関係ファイバー}
\label{sec:vi-setup-v28}

\subsection{三つの曲線と二つの基礎体}

$\mathcal R_3$ は、滑らかな種数三曲線 $C$ と非自明二捩れ
$\eta$ の対の粗モジュライとする。一般点では $C$ は非超楕円で、自己同型を持たない。
本稿では二つの関数体を区別する。
\[
 F=\mathbb C(Y_0(2)),\qquad K=\mathbb C(\mathcal R_3).
\]
$Y_0(2)$ は、楕円曲線とその位数二の部分群の粗モジュライである。
相対有理性は $F$ 上の五次元ファイバーの主張であり、
周期・指数は $K$ 上の曲線と曲面の主張である。

一般Prym対では $C/K$ は曲線として存在し、標準平面は
$\mathbb P(V^*)$、$V=H^0(C,\omega_C)$ という分裂した射影平面である。
一方、$\eta$ はまず
$\operatorname{Pic}^0_{C/K}(K)[2]$ の点として指定される。
線束自身の降下障害を
\[
 \alpha=\delta_C(\eta)\in\operatorname{Br}(K)[2]
\]
と書く。$\eta$ と $\eta^{\otimes2}\simeq\mathcal O_C$ を実現する
$\mu_2$-gerbe上では
$W=H^0(C,\omega_C\otimes\eta)$ は階数二の空間であり、
非自明な慣性は $W$ に $-1$ として作用する。
その射影化は、$K$ 上のPrym-canonical二次曲線 $D_\alpha$ へ降りる。

このgerbe上で基底を選べば、対称二次式表示を
\begin{equation}\label{eq:vi-symmetric-representation-v28}
 Q(x)=\begin{pmatrix}a(x)&b(x)\\ b(x)&c(x)\end{pmatrix},
 \qquad C=V(ac-b^2),
 \qquad N=\langle a,b,c\rangle
\end{equation}
と書ける。$\omega_C^2$ の切断の乗法から得る
$N\subset H^0(C,\omega_C^2)=\operatorname{Sym}^2V$ は $K$ に降りる。
これを二次曲線の\emph{net}と呼ぶ。
非退化二次曲線 $D\subset\mathbb P(N)$ を指定することは、
その双対 $D^\vee\subset\mathbb P(N^*)$ を指定することと同値である。
評価写像 $[x]\mapsto[a(x):b(x):c(x)]$ による $D^\vee$ の逆像が四次曲線 $C$ となる。
この双対化は基礎体上で定まるので、二次曲線に有理点を選ぶ必要はない。

$a=x^tA_0x$、$b=x^tA_1x$、$c=x^tA_2x$ と書く。
netのパラメータ平面にある判別式三次曲線は
\[
 \Delta_N=V\bigl(\det(\lambda_0A_0+\lambda_1A_1+\lambda_2A_2)\bigr)
 \subset\mathbb P(N)
\]
である。その対称一次表示の余核線束 $\theta_N$ は
\[
 0\to\mathcal O_{\mathbb P(N)}(-2)^3
 \xrightarrow{\lambda_0A_0+\lambda_1A_1+\lambda_2A_2}
 \mathcal O_{\mathbb P(N)}(-1)^3\to i_*\theta_N\to0
\]
により定まり、滑らかな場合には
$\theta_N^2\simeq\mathcal O_{\Delta_N}$、$H^0(\Delta_N,\theta_N)=0$ を満たす。

これとは別の、標準平面の双対平面内にある曲線を $E_N$ と書く。
双対平面上の普遍階数二部分束を $\mathcal S$ とし、
\[
 E_N=V(H_N),\qquad
 H_N=\det\bigl(N\otimes\mathcal O\longrightarrow
                       \operatorname{Sym}^2\mathcal S^*\bigr)
\]
と定める。これはnetの二次曲線に含まれる直線をパラメータ化する
\emph{Hermite曲線}である\cite[\S3.2.2, Exercise 3.12]{DolgachevCAG}。
一般netでは $E_N$ と $\Delta_N$ は滑らかであり、直線成分を含む退化二次曲線へ送る写像が
\[
 \pi_N:E_N\longrightarrow\Delta_N,
 \qquad (\pi_N)_*\mathcal O_{E_N}\simeq
                  \mathcal O_{\Delta_N}\oplus\theta_N
\]
という連結エタール二重被覆を与える。
従って $\operatorname{Pic}^0(E_N)$ は
$\operatorname{Pic}^0(\Delta_N)$ を $\langle\theta_N\rangle$ で割った楕円曲線である。
二つの種数一曲線の原点は指定しない。

この構成から得る古典的なnet写像は
\[
 \Psi:\mathcal R_3\dashrightarrow Y_0(2),
 \qquad (C,\eta)\longmapsto
 (\operatorname{Pic}^0(\Delta_N),\langle\theta_N\rangle)
\]
である。Dolgachev\cite[Remark 4.1]{DolgachevRationality2008}は、
これを一般netのモジュライへの写像として記述している。

\subsection{交代形式と関係代数からの幾何学的入力}

表示\eqref{eq:vi-symmetric-representation-v28}から
\[
 G=\begin{pmatrix}A_0&A_1\\A_1&A_2\end{pmatrix}
\]
を作り、$\det G\ne0$ の開部分を取る。
$[x]_\times v=x\times v$ とし、
$D_i=\operatorname{diag}([e_i]_\times,[e_i]_\times)$ と置く。
三つの交代形式を
\begin{equation}\label{eq:vi-omega-v28}
 \Omega_{ij}=D_i^tG^{-1}D_j-D_j^tG^{-1}D_i,\qquad
 \Omega(z)=z_1\Omega_{23}+z_2\Omega_{31}+z_3\Omega_{12}
\end{equation}
で定める。内在的には、これらは慣性weight一の階数六空間
\[
 T\simeq W^*\otimes\bigwedge^2V
\]
上の $\bigwedge^2V$ に値を持つ交代形式である。
一般の $E=E_N$ 上で $\Omega(z)$ の核は階数二のtwisted束 $\mathcal K$ となる。
その射影化が降りる $E$ 上の二次曲線束を $S\to E$ とする。
この構成の降下cocycleは第\ref{sec:prym-brauer-v27}節で確認する。

関係代数側の記号も固定する。
$L(x)=\sum x_iL_i$ は六行七列の一次式関係行列、$A=A^t\in\operatorname{GL}_7$ はpairingであり、
指定反対合の署名を $(2,4)$ とする。
$J=\operatorname{diag}(I_2,-I_4)$ として、課す全方程式は
\[
 L_iAL_j^tJ=JL_jAL_i^t\qquad(1\le i,j\le3)
\]
である。指定反対合を保持し、関係枠、二つの固有ブロックの基底、
共通の平面座標を除いた八次元有理商を $\mathfrak Z_3$ と書く。
その成分と良い開部分の構成は別稿\cite{PaperV}にある。

\begin{proposition}[本稿で用いる幾何学的事実]
\label{prop:vi-geometric-input-v28}
一般点について次が成立する。
\begin{enumerate}
\item 一般の種数三Prym対は対称表示\eqref{eq:vi-symmetric-representation-v28}を持つ。
表示の余核は $\omega_C\otimes\eta$ であり、選択差は二次元基底の合同変換と平面座標変更である。
従って三つの二次式の十八係数空間は $\mathcal R_3$ を支配する。
\item 同じ規約で
\begin{equation}\label{eq:vi-hermite-identity-v28}
 2\det(G)\operatorname{Pf}\Omega(z)=H_N(z).
\end{equation}
従って $E$ は交代行列のPfaffian曲線でもあり、
$\Omega(z)$ は $E$ 上で階数四を持つ。
\item 対応する写像
$\Phi_3:\mathfrak Z_3\dashrightarrow\mathcal R_3$ は支配的である。
その一般幾何学的ファイバーは $\mathbb P_E(\mathcal K)$ と双有理であり、
この同定はPrym表示の基底変更と可換である。
\end{enumerate}
\end{proposition}

対称表示と線束の対応には\cite{BeauvilleDeterminantal}を用いる。
関係成分からの支配性、式\eqref{eq:vi-hermite-identity-v28}、
第三項の同定の証明は\cite{PaperV}で与えた。
第三項に残る符号を明示すると、分裂体上で $T\oplus k^3$ の共通等方四平面の空間
\[
 X_Q=\{U\in\operatorname{Gr}(4,T\oplus k^3):
          (\Omega_{ij}\oplus\varepsilon_{ij})|_U=0\}
\]
を取り、$\tau(t,\ell)=(t,-\ell)$ の商を用いる。
良い開部分の $X_Q/\langle\tau\rangle$ が $\mathbb P_E(\mathcal K)$ と双有理になる。
三次元での外積二乗の持上げの二価性は、この $\tau$ がちょうど同一視する。
従ってこの幾何学的入力には、算術的降下の際にさらに選ぶべき平方根は含まれない。

\paragraph{計算に用いる一変数族.}
後で二次分裂体を具体的に求めるため、次の族を固定しておく。
\begin{equation}\label{eq:vi-hesse-example-v28}
 A_0=\begin{pmatrix}0&1&0\\1&0&0\\0&0&1\end{pmatrix},\quad
 A_1=\begin{pmatrix}1&0&0\\0&0&2\\0&2&0\end{pmatrix},\quad
 A_2=\begin{pmatrix}0&0&1\\0&t&0\\1&0&0\end{pmatrix}.
\end{equation}
このとき $\det G=t-4$ であり、
\[
 E_t:\quad 2tz_1^3+z_2^3+tz_3^3+(t-8)z_1z_2z_3=0.
\]
$t=1$ では滑らかである。式\eqref{eq:vi-hermite-identity-v28}は
\nolinkurl{hermite_pfaffian_v27.py} により普遍的な小行列式恒等式へ還元して検証される。

\paragraph{通常のPrym多様体との区別.}
上の補助種数一曲線の二次同種写像は、種数三Prym対に付随する
二次元Prym多様体の楕円因子を主張しない。
対称表示が定義される体上で、対応する非分岐二重被覆は
\[
 \widetilde C:\ a(x)=r^2,\quad b(x)=rs,\quad c(x)=s^2
          \subset\mathbb P^4
\]
という種数五のモデルを持つ。
Bruinの記述\cite[Theorem 5.1, case 4]{BruinPrymGenusThree}では、一般の場合、
そのPrym多様体は $y^2=-\det(A_0+2tA_1+t^2A_2)$ という種数二曲線のJacobianである。
本稿の $E\to\Delta_N$ は、平面内の退化二次曲線の二直線成分を選ぶエタール被覆である。

\section{net不変量を保つ相対有理性}
\label{sec:modular-fiber-rationality-v28}

種数三Prym対から二次曲線のnetを取り、その判別式三次曲線と非零二捩れを残す写像
\[
 \Psi:\mathcal R_3\dashrightarrow Y_0(2)
\]
自体は古典的である。Dolgachevの\cite[\S4, Remark 4.1]{DolgachevRationality2008}は、
この写像をnetのモジュライへの射影として記述している。
同論文はまた、一般netが非自明な安定化群を持つため、
粗モジュライ上のファイバーを直ちに射影空間束とみなせないことを指摘している。
ここではその安定化群を残したまま商を計算し、一般ファイバーの基礎体上の有理性を証明する。
別稿\cite{PaperV}のHermite--Pfaffian恒等式は、この古典的なnet写像を
関係代数のファイバーを記述するHermite曲線と結び付ける。

\begin{theorem}[net写像の相対双有理自明性]
\label{thm:modular-fiber-rationality-v28}
$F=\mathbb C(Y_0(2))$ とする。写像 $\Psi$ の一般ファイバーは
$F$ 上有理な五次元多様体である。従って $Y_0(2)$ の上で
\[
 \mathcal R_3\dashsim Y_0(2)\times\mathbb P^5
\]
という双有理同型が存在する。
この主張は、写像 $\Psi$ がその定義域全体で射影空間束となることを含まない。
\end{theorem}

\subsection{一般netの安定化群}

滑らかな判別式三次曲線を持つnet $N$ を取る。
その対称一次行列表示の余核を $\theta_N$ と書く。
対称最小分解の一意性から、一般の $N$ の安定化群は
\[
 \operatorname{Aut}(\Delta_N,\mathcal O_{\Delta_N}(1),\theta_N)
\]
と同型である。実際、三次曲線と余核線束の自己同型は最小分解へ持ち上がり、
対称性を用いて表示行列の合同変換となる。
逆に合同変換は余核線束を保つ。
安定化群からnetのパラメータ平面への作用は一般に忠実である。
その核の元はnetの全ての二次形式を同一のスカラー倍に移すが、
非退化な第一形式を単位行列、第二形式を相異なる固有値を持つ対角行列にした後、
第三形式の一般の非対角成分を用いると、その元は射影的に恒等となる。

代数閉包上で三次曲線の変曲点を原点に選べば
$\mathcal O_{\Delta_N}(1)\simeq\mathcal O_{\Delta_N}(3\cdot0)$ となる。
一般の $j$ 不変量に対する安定化群は
\[
 G\simeq\Delta_N[3]\rtimes\{\pm1\}
       \simeq (\mathbb Z/3)^2\rtimes\mathbb Z/2
\]
であり、位数は十八である。平行移動は $\operatorname{Pic}^0$ に自明に作用し、
$[-1]$ は二捩れを保つため、$\theta_N$ の指定はこの群を縮めない。
Hesse座標では、この射影作用は
\[
 s(x,y,z)=(y,z,x),\qquad
 t(x,y,z)=(x,\zeta y,\zeta^2z),\qquad
 r(x,y,z)=(x,z,y)
\]
により生成される。ここで $\zeta^3=1$、$\zeta\ne1$ である。
Hesse標準形とこの作用については\cite[\S2, \S4]{ArtebaniDolgachevHesse}を参照されたい。

\subsection{位数十八の射影群による二次曲線空間の商}

次の補題は、安定化群が基礎体上で定数群でない場合にも適用できる。

\begin{lemma}[安定化群の降下を含む有理性]
\label{lem:hesse-eighteen-conics-v28}
$k$ を標数零の体、$V$ を三次元の $k$ ベクトル空間とする。
有限エタール部分群スキーム $G\subset\operatorname{PGL}(V)$ が、
代数閉包上で上の位数十八の射影群と共役であると仮定する。
このとき
\[
 \mathbb P(V)/G\simeq\mathbb P_k(1,1,2),
 \qquad
 \mathbb P(\operatorname{Sym}^2V^*)/G\dashsim_k\mathbb P_k^5.
\]
第二の商は、第一の商の稠密開部分上の階数三ベクトル束と双有理である。
\end{lemma}

\begin{proof}
$\widetilde G\subset\operatorname{SL}(V)$ を $G$ の逆像とする。
幾何学的には位数五十四であり、唯一のSylow三部分群で割ると位数二となる。
従って符号指標
\[
 \chi:\widetilde G\longrightarrow\{\pm1\}
\]
は $k$ 上で定まる。表現をこの指標でねじり、
\[
 V'=V\otimes\chi,\qquad
 \widehat G=\{\chi(g)g:g\in\widetilde G\}\subset\operatorname{GL}(V)
\]
と置く。このねじりは射影作用を変えない。
幾何学的には $\widehat G$ は、成分の積が一である三乗根の対角行列の群 $D$ と
三変数の置換群の半直積
\[
 \widehat G_{\overline k}=D\rtimes\mathfrak S_3
\]
となる。これは位数五十四の単項行列群である。

まず不変式環を計算する。$D$ に不変な単項式 $x^ay^bz^c$ は
$a\equiv b\equiv c\pmod3$ を満たす。
したがって $(xyz)^e$、$e=0,1,2$ を除けば $x^3,y^3,z^3$ の単項式である。
さらに置換群で不変なものを取ると、基本対称式の定理から
\[
 \overline k[V']^{\widehat G}
 =\overline k[p,q,r],\qquad
 p=x^3+y^3+z^3,\quad q=xyz,\quad
 r=x^3y^3+y^3z^3+z^3x^3
\]
を得る。次数はそれぞれ三、三、六である。
ここで第三基本対称式は $q^3$ なので別の生成元を要しない。

この表示は $k$ 上にも降りる。
有限群の不変式は標数零で体の拡大と可換であり、次数三の不変式空間は $k$ 上二次元である。
その基底 $p_0,q_0$ を取り、次数六の不変式空間内で
$\operatorname{Sym}^2\langle p_0,q_0\rangle$ の補空間から非零元 $r_0$ を取る。
代数閉包上で生成と代数的独立性が成立するので、忠実平坦降下から
\[
 k[V']^{\widehat G}=k[p_0,q_0,r_0].
\]
従って $\mathbb P(V)/G$ は分裂した重み付き射影平面
$\mathbb P_k(3,3,6)\simeq\mathbb P_k(1,1,2)$ である。

次に二次曲線の表現を見る。
上の単項行列群に対し、二次式の二つの部分空間
\[
 \langle x^2,y^2,z^2\rangle,
 \qquad \langle yz,zx,xy\rangle
\]
は、ともに三次元表現 $V'$ と同型である。
これは対角行列と置換行列で直接確認できる。
よって $\overline k$ 上で
\[
 \operatorname{Sym}^2V^*\simeq V'\oplus V'
\]
となる。基礎体への降下では、二つの部分空間自体をそれぞれ降ろす必要はない。
代わりに
\[
 M=\operatorname{Hom}_{\widetilde G}
       (V',\operatorname{Sym}^2V^*)
\]
を取ると、$M$ は $k$ 上二次元であり、評価写像
\[
 V'\otimes_kM\longrightarrow\operatorname{Sym}^2V^*
\]
は代数閉包上で同型なので $k$ 上でも同型である。
ここで $\chi^2=1$ により、二次式の表現は $V$ を $V'$ に変えても同じである。
$M$ の $k$ 基底を選ぶと、求める射影空間は
$\mathbb P(V'\oplus V')$ と $G$ 同変に同型になる。

第一成分への射影
\[
 \mathbb P(V'\oplus V')\dashrightarrow\mathbb P(V')
\]
の、第一成分が非零である開部分は
$\operatorname{Hom}(\mathcal O_{\mathbb P(V')}(-1),V')$ の全空間である。
中心 $\mu_3\subset\widetilde G$ はこのHom束では相殺されるので、
これは真の $G$ 同変階数三ベクトル束である。
$G$ が自由に作用する底の稠密開部分へ制限して商を取れば、
有限エタール降下により、$\mathbb P(V')/G$ の開部分上の階数三ベクトル束を得る。
従って商の関数体は $k(\mathbb P(V')/G)$ に三つの独立変数を加えたものとなる。
第一項の重み付き射影平面は $k$ 有理なので、第二項も $k$ 有理である。
\end{proof}

\subsection{モジュラー関数体上のnet代表と定理の証明}

一般netの代表をモジュラー関数体そのものの上で選べることも明示する。
$u$ を変数とし、
\[
 N_u=\langle y^2+2xz,\ 2yz,\ -x^2+2xz-2y^2+uz^2\rangle
 \subset\operatorname{Sym}^2\mathbb C(u)^3
\]
とする。この族はnetの古典的な正規形の一つである
\cite[\S3.2, \S5]{DomokosFeherRimanyi2013}。
判別式の対称行列は
\[
 A(\lambda)=
 \begin{pmatrix}
 -\lambda_2&0&\lambda_0+\lambda_2\\
 0&\lambda_0-2\lambda_2&\lambda_1\\
 \lambda_0+\lambda_2&\lambda_1&u\lambda_2
 \end{pmatrix}.
\]
$\lambda_2=1$、$X=\lambda_0-2$、$Y=\lambda_1$ と置くと、
その行列式の零点集合は
\[
 \Delta_u:\quad Y^2=X(X^2+6X+9+u)
\]
となる。$u\ne0,-9$ では滑らかであり、$\mathbb C(u)$ 上の非零有理二捩れ点は
$T=(0,0)$ だけである。
従って対称表示から得る非零二捩れ $\theta_{N_u}$ は $T$ に一致する。
その $j$ 不変量は
\[
 j(\Delta_u)=\frac{1728(u-3)^3}{u(u+9)^2}.
\]
この有理関数の次数は三であり、$Y_0(2)$ から $j$ 直線への忘却写像も次数三である。
したがって $u$ は $Y_0(2)$ の双有理座標であり、$N_u$ は
$F=\mathbb C(Y_0(2))$ 上の一般netの代表を与える。

\begin{proof}[定理\ref{thm:modular-fiber-rationality-v28}の証明]
一般のPrym対の対称二行二列二次表示は、net $N$ の中の非退化二次曲線を指定することと同値である。
ここで二次曲線は $\mathbb P(N)$ 内に取り、その双対二次曲線を
評価写像 $\mathbb P^2\to\mathbb P(N^*)$ で引き戻す。
二次曲線の双対化は基礎体上で定義され、射影基底変更と可換である。
二次曲線を非特異な標準形へ移した後に $Q=\left(\begin{smallmatrix}a&b\\b&c\end{smallmatrix}\right)$
を取れば、元の四次曲線は $ac-b^2=0$ として復元される。
ここでは粗モジュライを扱うため、基礎体上でその二次曲線に有理点があることは要求しない。
このnetと二次曲線による表示は\cite[Remark 4.1]{DolgachevRationality2008}の記述と一致する。

一般netの幾何学的同型類は $(\Delta_N,\theta_N)$ により決まる。
このことは対称一次行列表示の一意性から従う。
また代数閉包上では、種数一曲線の次数三の埋込み線束の選択差は平行移動で吸収される。
よってnetの粗モジュライは $Y_0(2)$ と双有理である。

上で得た $F$ 上の代表 $N_u$ を固定し、そのパラメータ空間を $V=N_u$ とする。
一般ファイバーは、一般netの軌道上の二次曲線束を射影座標変更で割ったものであり、
その関数体は
\[
 F\bigl(\mathbb P(\operatorname{Sym}^2V^*)\bigr)^{\operatorname{Aut}(N_u)}
\]
である。これは安定化群を保った商であって、単に二次曲線空間を底へ降ろしたものではない。
その群は前節の位数十八の群の $F$ 形なので、
補題\ref{lem:hesse-eighteen-conics-v28}によりこの関数体は
$F$ の五変数純超越拡大である。
\end{proof}

\paragraph{算術的ファイバーとの両立.}
この相対有理性は、$K=\mathbb C(\mathcal R_3)$ 上の一般Hermite曲線 $E$ や
Prym-canonical二次曲線 $D_\alpha$ の分裂を意味しない。
一般netの代表を固定した二次曲線空間から粗モジュライへ移る際には、
上の有限安定化群による降下が残る。
従って $E$ の周期・指数三、$D_\alpha$ の指数二、
および $E\times D_\alpha$ の指数六という主張とは矛盾しない。

\paragraph{再現入口.}
\nolinkurl{modular_fiber_rationality_v28.py} は、位数五十四の単項行列群と
位数十八の射影像、二次式表現の二重化、不変式の次数、
正規形netの行列式と $j$ 不変量を厳密に検証する。
不変式環の全次数にわたる生成、群スキーム上の降下、およびベクトル束の降下は、
上の証明による。

\section{一般Hermite曲線の周期と指数}
\label{subsec:hermite-index-three-v27}

ここでは $K=\mathbb C(\mathcal R_3)$ とし、一般Prym対に付随するHermite曲線を $E/K$ と書く。
曲線の指数は、その閉点の剰余体次数の最大公約数である。
種数一曲線の周期は、$[E]\in H^1(K,\operatorname{Pic}^0(E))$ というtorsor classの位数とする。
一般の曲線が標数零代数閉体上で種数一であることは、$K$ 上で原点を持つことを意味しない。

\begin{theorem}[補助曲線の指数三]
\label{thm:hermite-index-three-v27}
一般補助曲線 $E/\mathbb C(\mathcal R_3)$ の周期と指数はともに三である。
特に $E(K)=\varnothing$ であり、原点を有理的に選ぶことはできない。
\end{theorem}

\begin{proof}
まず上界を確認する。一般の種数三曲線は非超楕円で自己同型を持たず、
その標準平面は $\mathbb P(H^0(C,K_C)^*)$ という分裂した射影平面である。
net $N\subset H^0(C,K_C^2)$ は対称Prym表示の選択に依存せず、
Hermite曲線はこの標準平面の双対平面内の三次曲線として $K$ 上へ降りる。
従って一般の $K$ 上の直線との交点は次数三の零サイクルを与え、
$\operatorname{ind}(E)\mid3$ である。

下界には一つの支配的な族を用いる。
$r$ を超越元とし、離散付値パラメータ $s$ を使って
\[
 M_s(u,v,w)=
 \begin{pmatrix}
 ru&sw&sv\\
 sw&rsv&su\\
 sv&su&rs^2w
 \end{pmatrix}
\]
を取る。これは二次曲線のnetの対称表示であり、
\[
 \det M_s=-rs^2\left(
 u^3+sv^3+s^2w^3-\frac{r^3+2}{r}suvw\right).
\]
対応するHermite曲線は
\begin{equation}\label{eq:index-three-hermite-v27}
 E_s:\quad
 r^2(s^2z_1^3+sz_2^3+z_3^3)+(r^3-4)sz_1z_2z_3=0.
\end{equation}
これらの行列式の恒等式は展開によって得られる。
$t^3=s$ を添加し、判別式曲線では
$(U,V,W)=(u,tv,t^2w)$ と置くと
\[
 M_s(U,V/t,W/t^2)
 =\operatorname{diag}(1,t,t^2)
   M_1(U,V,W)
   \operatorname{diag}(1,t,t^2)
\]
となる。Hermite曲線では $(Z_1,Z_2,Z_3)=(t^2z_1,tz_2,z_3)$ を使う。
二つの三次曲線の滑らかさを与える因子はそれぞれ
\[
 (r^3-1)^2(r^3+8),\qquad (r^3-1)(r^3+8)^2
\]
なので、$r$ が超越元なら両方とも滑らかである。

このnet族がPrymモジュライを支配することも確認する。
$s=1$ とし、$M_1$ の三係数行列が定める順序付き二次式を $(q_0,q_1,q_2)$ とする。
平面座標の $g\in\operatorname{GL}_3$ とnet基底の $h\in\operatorname{GL}_3$ を使い、
\[
 (r,g,h)\longmapsto h\,(q_0(gx),q_1(gx),q_2(gx))^t
\]
を考える。これは十九変数から十八個の二次式係数への写像である。
$r=2,g=h=I$ において、その接行列の十八次小行列式の一つは
$-462422016$ となる。従ってこの写像は支配的である。
$s$ も変数とした同じ小行列式は零多項式ではないので、$\mathbb C(s)$ 上の族も支配的である。
一般の順序付き二次式を $Q=\left(\begin{smallmatrix}q_0&q_1\\q_1&q_2\end{smallmatrix}\right)$
へ入れれば、命題\ref{prop:vi-geometric-input-v28}の対称表示を通じて $\mathcal R_3$ を支配する。
必要な滑らかさ・非退化性はすべて稠密開条件として満たされる。

$g,h$ の成分を独立な超越元 $\mathbf a$ とし、
\[
 F=\mathbb C(r,\mathbf a)((s))
\]
と置く。上の支配性から $K\hookrightarrow F$ を得る。
基底変更 $h$ はnetそのものを変えず、$g$ はHermite曲線を射影同型に移すので、
$E_F$ は式\eqref{eq:index-three-hermite-v27}の $E_s$ と同型である。

有限拡大 $L/F$ の次数が三で割れないと仮定する。
標数零の完備離散付値体の基本等式から、分岐指数 $e=v_L(s)$ も三で割れない。
$r$ と $r^3-4$ は単元である。もし $E_s(L)$ に点があり、三座標が全て非零なら、
それらの付値を $a,b,c$ として、三つの立方項の付値は
\[
 2e+3a,\qquad e+3b,\qquad 3c
\]
となる。これらは三を法として相異なり、最小値は一つだけである。
混合項の付値は
\[
 e+a+b+c=\frac{(2e+3a)+(e+3b)+3c}{3}
\]
なので、その最小値より真に大きい。従って最小付値の項を打ち消すことができず、矛盾する。
ある座標が零の場合も、残る立方項の付値が三を法として異なるため同じ矛盾を得る。
よって $E_s$ の全閉点の次数は三の倍数である。
一方、$t^3=s$ 上には
$[z_1:z_2:z_3]=[t^{-2}:-t^{-1}:0]$ という点があるから、
$\operatorname{ind}(E_F)=3$ である。

基礎体を拡大すると指数は元の指数を割るので、
$3=\operatorname{ind}(E_F)\mid\operatorname{ind}(E)$ となる。
上界と合わせて指数は三である。
種数一曲線の周期は指数を割り、周期一ならそのJacobianの自明torsorとして有理点を持つ。
従って周期も三である。
\end{proof}

\section{一般PrymファイバーのBrauer障害}
\label{sec:prym-brauer-v27}

命題\ref{prop:vi-geometric-input-v28}の幾何学的ファイバーの記述を、粗モジュライの関数体上へ降ろす。
ここでは、曲線上の線束の同型類が基礎体上で定まることと、
その線束自身が降下することを区別する。
この違いが、本族の一般ファイバーに実際のねじれを生じる。

\subsection{Prym-canonical pencilの降下障害}
一般の四次曲線に自己同型がない開部分を取り、
\[
 K=\mathbb C(\mathcal R_3),\qquad
 C/K,\qquad
 0\ne\eta\in\operatorname{Pic}^0_{C/K}(K)[2]
\]
を一般Prym対とする。
曲線 $C$ は $K$ 上に存在するが、$\eta$ はまずPicardスキームの点として指定される。
線束の降下障害を
\[
 \alpha=\delta_C(\eta)\in\operatorname{Br}(K)[2]
\]
と書く。
ここで $\delta_C$ は、幾何学的線束のGalois降下同型の三重積に現れる
スカラー二次cocycleを取る写像である。
Brauer群と射影線形群の降下cocycleの対応には
\cite[Tags 03R1, 0A2J]{StacksBrauerV27}を用いる。

$V=H^0(C,\omega_C)$ は三次元の $K$ ベクトル空間である。
一方、
\[
 W=H^0(C_{\overline K},\omega_C\otimes\eta)
\]
は二次元であり、その降下行列も $\alpha$ を表すスカラー二次cocycleを持つ。
従って射影化は $K$ 上のBrauer--Severi二次曲線へ降りる。
これを $D_\alpha$ と書く。幾何学的には、これはPrym-canonical pencilの標的である。

\begin{theorem}[一般ファイバーの非自明なねじれ]
\label{thm:prym-brauer-v27}
一般Prym対上で次が成立する。
\begin{enumerate}
\item 補助Pfaffian三次曲線は、$K$ 上の分裂した射影平面内の
滑らかな曲線 $E/K$ へ降りる。
$\Phi_3$ の一般ファイバーは、$E$ 上のBrauer--Severi二次曲線束
$\pi:S\to E$ と $K$ 双有理である。そのBrauer classは $\alpha|_E$ であり、
\[
 S\dashsim_K E\times_KD_\alpha
\]
となる。
\item $\alpha$ は零でなく、位数も指数も二である。
さらに $\alpha|_{K(E)}\ne0$ であり、$S\to E$ は一般点でも分裂しない。
従って $S$ は $E$ 上の真の階数二ベクトル束の射影化とはならない。
一方、$\alpha$ は $K(S)=\mathbb C(\mathfrak Z_3)$ 上では零になる。
\item 一般ファイバーの滑らかな射影モデルの指数は
\[
 \operatorname{ind}(S)=2\operatorname{ind}(E)=6.
\]
さらに $S$ の任意の非空開部分は次数六の閉点を持つ。
従って、一般ファイバーを正則な開部分に取った有理商の多重切断の最小次数は六であり、
全ての多重切断の次数は六の倍数である。
\end{enumerate}
\end{theorem}

ここで指数とは閉点の次数全体の最大公約数である。
切断の主張は滑らかなモデル、またはその正則な開部分に関するものである。
特異な境界への切断まで任意の双有理モデルで禁止するという意味ではない。

\subsection{核束の降下cocycle}
\begin{proof}[定理\ref{thm:prym-brauer-v27}の第一項]
$\eta$ を線束として代表し、$\eta^{\otimes2}\simeq\mathcal O_C$ の同型を選ぶ
$\mu_2$-gerbe上で構成を行う。
その慣性の非自明元は $W$ に $-1$ として作用する。
標準平面四次曲線に対して
$\operatorname{Sym}^2V\simeq H^0(C,\omega_C^{\otimes2})$ なので、
切断の乗法は
\[
 \operatorname{Sym}^2W\longrightarrow\operatorname{Sym}^2V
\]
を与える。これは二次対称行列 $Q$ の座標を外した記述である。
この規約での対称表示の余核は、正確には
$(\omega_C\otimes\eta)\otimes\det(W)^{-1}$ である。
$\det(W)$ には慣性が自明に作用するので、これは $K$ 上に降る一次元線である。
以下ではこの基礎体上の線を自明化する。
$Q$ は $\operatorname{Sym}^2W^*\otimes\operatorname{Sym}^2V$ の元と解釈する。
$H=W\otimes V^*$ の上に係数Gram形式 $G$ が得られ、一般に非退化である。

一次直交syzygyの空間を $T$ とする。
その元 $b$ は $V^*\to H$ という一次写像であり、
全ての $w\in W$ と $x\in V^*$ に対して
$G(w\otimes x,b(x))=0$ を満たす。
$G$ で値を双対化すると、この条件は二つの $V$ 成分の反対称性になるので、
\[
 T\simeq W^*\otimes\bigwedge^2V,\qquad \dim T=6.
\]
特に慣性の非自明元は $T$ にも $-1$ として作用する。

$b,c\in T$ に対し、
\[
 \Omega(b,c)(x,y)=G(b(x),c(y))-G(b(y),c(x))
\]
と定める。これは基底に依存しない写像
$\bigwedge^2T\to\bigwedge^2V$ であり、
式\eqref{eq:vi-omega-v28}の三つの交代形式を与える。
形式のパラメータ平面は
\[
 \mathbb P\bigl((\bigwedge^2V)^*\bigr)\simeq\mathbb P(V)
\]
である。従って標準平面 $\mathbb P(V^*)$ の双対として $K$ 上分裂している。
Pfaffianの零点集合は慣性に不変なので $E/K$ へ降りる。

$E$ 上の交代写像の核 $\mathcal K$ は階数二であり、
慣性はその上で $-1$ として作用する。
従って $\mathbb P(\mathcal K)$ は $E$ 上の二次曲線束 $S$ へ降り、
その降下cocycleは $\alpha|_E$ である。
同時に、命題\ref{prop:vi-geometric-input-v28}の $X_Q/\langle\tau\rangle$ との双有理対応は
これらの基底変更と可換である。
外積二乗の持上げの二価性は既に $\tau$ の商で同一視されているので、
$K$ 上で新たな平方根を選ぶ必要はない。
よって $S$ は実際の一般ファイバーの $K$ 双有理モデルとなる。

$K(E)$ 上では、$S$ の一般二次曲線と $D_\alpha$ の基底拡大は
同じBrauer classを持つ。
同じ次数と同じBrauer classを持つ中心単純代数は同型なので、
この二つの二次曲線は同型である。
これが $S\dashsim_K E\times_KD_\alpha$ を与える。
ここでは一般点上の同型を用いており、$E$ 全体の上での束の同型までは主張しない。
\end{proof}

\subsection{一般Brauer classの非消滅}
\begin{lemma}[非分裂二次曲線からの支配族]
\label{lem:prym-brauer-nonzero-v27}
$\delta_C(\eta)$ は $\mathbb C(\mathcal R_3)$ 上で零でない。
\end{lemma}

\begin{proof}
$F=\mathbb C(s,t)$ とし、
\[
 D:\quad Z^2-sX^2-tY^2=0\subset\mathbb P^2_F
\]
を取る。この二次曲線は、任意の独立な変数組 $\mathbf u$ を加えた
$F(\mathbf u)$ 上でも分裂しない。
実際、$s$ 進付値で射影座標の最小付値を零にすると、
剰余体 $\mathbb C(t,\mathbf u)$ 上で $Z_0^2=tY_0^2$ となる。
$t$ は平方でないので $Z_0=Y_0=0$、従って $X_0\ne0$ である。
元の方程式では $sX^2$ だけが付値一、他の二項が付値二以上となり、矛盾する。

十八個の独立係数を持つ三つの二次式 $a,b,c$ を取り、
\[
 f:\mathbb P^2\longrightarrow\mathbb P^2,
 \qquad [x]\longmapsto[a(x):b(x):c(x)],
 \qquad C=f^{-1}(D),\quad\rho=f|_C
\]
とする。一般の係数では基点がなく、$C$ は滑らかな平面四次曲線である。
この開条件の非空性は、例えば $a=x_1^2,b=x_2^2,c=x_3^2$ からも分かる。

$\ell\in\operatorname{Pic}_{D/F}(F)$ を幾何学的次数一のclassとし、
\[
 M=\rho^*\ell,
 \qquad \eta=M\otimes\omega_C^{-1}
\]
と置く。$\ell^{\otimes2}=\mathcal O_D(1)$ と
$\rho^*\mathcal O_D(1)=\mathcal O_C(2)=\omega_C^{\otimes2}$ より、
$\eta^{\otimes2}=\mathcal O_C$ である。
代数閉包上で $D$ をVeronese二次曲線にすれば、
$M$ は対称二行二列二次式行列の余核になる。
その分解
\[
 0\longrightarrow\mathcal O(-2)^2
 \longrightarrow\mathcal O^2\longrightarrow i_*M\longrightarrow0
\]
から $h^0(M)=2$ である。
$h^0(\omega_C)=3$ なので $\eta$ は非自明である。

降下障害の自然性から
\[
 \delta_C(\eta)=\delta_D(\ell)=[D]\ne0.
\]
代数閉包上では、三つの二次式の十八係数は
対称二行二列二次式行列 $Q$ の全係数を動かす。
命題\ref{prop:vi-geometric-input-v28}の対称表示の支配性により、この族は
$\mathcal R_3$ を支配する。
従って一般の $\alpha$ が零なら、この支配族上でも零になり、矛盾する。
\end{proof}

この非消滅は明示的なcocycleでも確認できる。
$r^2=s$ と置いた二次拡大上で
\[
 Q=\begin{pmatrix}c+ra&b\\ b&(c-ra)/t\end{pmatrix},
 \qquad T=\begin{pmatrix}0&t\\1&0\end{pmatrix}
\]
とすると、$\sigma(r)=-r$ に対して
\[
 T\sigma(Q)T^t=tQ,\qquad T\sigma(Q)=QT^t,
 \qquad T\sigma(T)=tI_2.
\]
従って余核線束の半線形降下作用の二乗は $t$ であり、
その障害は四元数代数 $(s,t)$ のclassである。
付属の \nolinkurl{brauer_conic_cocycle_v27.py} は、
この三恒等式と行列式の式を有理関数体上で照合する。
支配性と非消滅の根拠は、上の証明そのものである。

\subsection{三次曲線への制限と切断の不存在}
\begin{proof}[定理\ref{thm:prym-brauer-v27}の残りの主張]
補題\ref{lem:prym-brauer-nonzero-v27}と二次曲線表示により、
$\alpha$ の位数と指数は二である。
$E$ は $K$ 上の平面三次曲線なので、直線との交点として次数三の零cycleを持つ。
もし $\alpha|_{K(E)}=0$ なら、$D_\alpha$ の $K(E)$ 点から
有理写像 $E\dashrightarrow D_\alpha$ を得る。
$E$ が滑らかな曲線で $D_\alpha$ が固有なので、この写像は全体に延びる。
次数三の零cycleを押し出すと、$D_\alpha$ 上に次数三の $K$ 因子ができる。
その線束に $\omega_{D_\alpha}$ を掛けると次数一となり、
種数零のRiemann--Rochから非零切断を持つ。
従って次数一の有効因子、すなわち $K$ 点が存在し、非分裂性に矛盾する。
よって $\alpha|_{K(E)}\ne0$ である。
$S$ の一般点はその二次曲線上の点を与えるので、$\alpha|_{K(S)}=0$ となる。

$E\times D_\alpha$ の両射影により、その指数は
$\operatorname{ind}(E)$ と $2$ の最小公倍数で割り切れる。
零cycleの外積により、その指数は $2\operatorname{ind}(E)$ を割る。
$\operatorname{ind}(E)\mid3$ なので両者は等しく、
\[
 \operatorname{ind}(E\times D_\alpha)=2\operatorname{ind}(E).
\]
滑らかな射影曲面の点吹上げは指数を変えないため、
不定点の解消により同じ式が $S$ と全ての滑らかな射影双有理モデルに成立する。
定理\ref{thm:hermite-index-three-v27}で証明した
$\operatorname{ind}(E)=3$ を用いると、$\operatorname{ind}(S)=6$ を得る。

指数が六であることだけでは、次数六の閉点の存在までは従わないので、これを別に示す。
$U\subset E\times D_\alpha$ を任意の非空開部分とする。
射影の像は $E$ の非空開部分 $E_0$ であり、その補集合は有限である。
$K$ は無限体なので、$E\setminus E_0$ を避けて $E$ と横断的に交わる $K$ 直線を選べる。
その交点因子は次数三である。
$\operatorname{ind}(E)=3$ なので、これは一つの次数三の閉点 $x$ からなる。
$L=\kappa(x)$ と置くと、$D_{\alpha,L}$ はなお非分裂である。
もし分裂すれば $D_\alpha$ に次数三の零cycleができ、上と同じ種数零の議論に反する。

$U_x$ は $D_{\alpha,L}$ の非空開部分である。
$L$ 上の直線を、その補集合を避けて二次曲線と横断的に交わるよう選ぶ。
非分裂性により、その次数二の交点は一つの閉点 $y$ である。
従って
\[
 [\kappa(y):K]=[\kappa(y):L][L:K]=2\cdot3=6,
 \qquad y\in U.
\]
双有理同型の定義される開部分へ制限すれば、$S$ の任意の非空開部分にもこの構成を移せる。

正則な一般ファイバーの任意の次数 $d$ の閉点からは、
Lang--Nishimuraの補題により、その剰余体上の $S$ の点を得る。
この曲面の場合には、固有な $S$ への有理写像の不定点を点吹上げで解消し、
例外曲線が剰余体上の射影直線であることからも同じ結論が従う。
従って $6\mid d$ である。
一方、上で構成した次数六の閉点の閉包は次数六の多重切断を与える。
これで最小次数が六であることが分かる。
\end{proof}

\paragraph{未決の境界.}
この結果は、幾何学的ファイバーの射影直線束が一般関数体上では実際にねじれることを決定する。
ねじれを保持した良いコンパクト化、
標識を忘れた関係代数のモジュライへの降下は、ここでは決定していない。
また、一般ファイバーの記述から全八次元空間の有理性を結論しない。

\section{Pl\"ucker二次曲線と内在的な二次分裂体}
\label{sec:intrinsic-quadratic-splitting-v28}

一般Prym対のBrauer障害には、補助スクロール自身から定まる二次分裂体がある。
この体は、三つの交代形式の共通等方三平面を区別する体であり、
二十次元から十八次元への一つの線形写像と、二変数二次式から計算できる。
以下では引き続き $K=\mathbb C(\mathcal R_3)$ とし、
$\alpha=\delta_C(\eta)\ne0$ とする。
前節の $\mu_2$-gerbe上の階数六空間を $T$、
三次元の通常の $K$ 空間を $B=\bigwedge^2V$ と書く。
交代形式は一つの写像 $\Omega:\bigwedge^2T\to B$ にまとめる。

\begin{theorem}[内在的な二次分裂体]
\label{thm:intrinsic-quadratic-splitting-v28}
一般点上で次が成立する。
\begin{enumerate}
\item 縮約写像
\[
 c_\Omega:\bigwedge^3T\longrightarrow T\otimes B,\qquad
 u\wedge v\wedge w\longmapsto
 w\otimes\Omega(u,v)-v\otimes\Omega(u,w)+u\otimes\Omega(v,w)
\]
は階数十八であり、$\Pi=\ker c_\Omega$ は階数二である。
射影化 $D_\Pi=\mathbb P(\Pi)$ は $K$ 上の二次曲線へ降り、
そのBrauer classは $\alpha$ である。
従って $D_\Pi$ はPrym-canonical pencilの二次曲線 $D_\alpha$ と $K$ 同型である。
この同型の選択は主張に含めない。
\item 共通等方三平面のスキーム $Y$ は、Pl\"ucker埋込みによって
\[
 Y=\operatorname{Gr}(3,T)\cap\mathbb P(\Pi)\subset D_\Pi
\]
と書ける。これは幾何学的次数二の被約因子であり、
$Y=\operatorname{Spec}L$ は非自明な二次体拡大 $L/K$ を定める。
この体上で $\alpha_L=0$ となる。
\item $L$ 上では $T_L=U_+\oplus U_-$ と分解し、
$U_\pm$ は二つの共通等方三平面である。
補助三次曲線上の核束も
\[
 \mathcal K_L=\mathcal M_+\oplus\mathcal M_-,\qquad
 \deg\mathcal M_+=\deg\mathcal M_-=-3
\]
と分かれる。非自明なGalois元は、元のtwisted束の降下、従ってその射影化の上で二つの項を交換する。
\item 二次拡大後も $E_L$ の周期と指数は三であり、
$S_L\dashsim_L E_L\times\mathbb P^1_L$ かつ
$\operatorname{ind}(S_L)=3$ である。
\end{enumerate}
\end{theorem}

\begin{proof}
まず幾何学的な一般点における線形代数を確認する。
十八係数のGram空間上で $G=\operatorname{diag}(I_3,I_3)$ を取る。
$T=A\oplus A$ の基底を $e_1,e_2,e_3,f_1,f_2,f_3$ とし、
交代形式の符号を共通に変更すると
\[
 \Omega((v,w),(v',w'))=v\times v'+w\times w'
\]
となる。
$p_j$ を
\[
 (e_1+xf_1)\wedge(e_2+xf_2)\wedge(e_3+xf_3)
   =p_0+xp_1+x^2p_2+x^3p_3
\]
の係数とする。
$\bigwedge^3T$ を $e$ と $f$ の個数で四つに分けて縮約を計算すると、
\[
 \ker c_\Omega=\langle p_0-p_2,\ p_1-p_3\rangle
\]
を得る。具体的には、$\bigwedge^3A$ と
$A\otimes\bigwedge^2A$ から $A\otimes B$ への条件が九本あり、
その核は $p_0-p_2$ の直線である。
もう一方の九本は $p_1-p_3$ の直線を残す。
従って縮約の階数は十八である。

この射影直線上のPl\"ucker二次関係式は、全て
\[
 a^2+b^2=0,\qquad
 [a:b]\longmapsto[a(p_0-p_2)+b(p_1-p_3)]
\]
の倍数であり、少なくとも一つは非零の倍数である。
二点は $U_\pm=\{(v,\pm\sqrt{-1}v):v\in A\}$ を表し、
被約かつ互いに相補的である。
また、任意の三次元部分空間の非零Pl\"uckerベクトル $\xi$ に対して、
$c_\Omega(\xi)=0$ は、その部分空間が三つの形式に共通等方的であることと同値である。
基底 $u,v,w$ を取れば、縮約の三つの係数が
$\Omega(u,v),\Omega(u,w),\Omega(v,w)$ であることから従う。
これは相対Grassmann多様体上でも同じ零点条件を与える。

階数十八は開条件である。
上の点では相対Grassmann多様体の零点スキームは、期待次元零の被約な二点である。
従ってその近傍でエタールであり、射影性により余分な点が出現しない開近傍を取れる。
この開部分では $Y$ は次数二の有限エタールスキームで、二つの平面は相補的である。
Gram空間は既約であり、滑らかなPrym四次曲線を与える開部分は非空なので、
二つの開条件を同時に満たす一般点を取れる。
Gram表示からPrymモジュライへの支配性により、一般Prym対にも適用できる。

次に降下を考える。
慣性の非自明元は $T$ と $\bigwedge^3T$ にともに $-1$ として作用する。
$c_\Omega$ は慣性と可換なので、$\Pi$ もweight一である。
従って $\mathbb P(\Pi)$ の降下cocycleは $\alpha$ を表す。
次数二の中心単純代数はBrauer classによって同型まで決まるので、
$D_\Pi\simeq_KD_\alpha$ となる。
Pl\"ucker埋込みもscalar慣性を消すため降下し、$Y\subset D_\Pi$ が得られる。
幾何学的には $\mathbb P^1$ 上の次数二因子なので、これは反標準因子の類を持つ。

$\alpha\ne0$ だから $D_\Pi$ は $K$ 点を持たない。
従って次数二エタールスキーム $Y$ は二つの $K$ 点には分かれず、
二次体 $L$ のスペクトルである。
そのtautologicalな点は $D_\Pi(L)$ の点でもあるので、$\alpha_L=0$ となる。
同じ結論は、tautologicalなweight一・階数三部分空間の行列式が
weight一の線になることからも得られる。

$L$ 上でgerbeを中立化すると、$T_L=U_+\oplus U_-$ は通常の空間の直和である。
この分解に対して
\[
 \Omega(z)=\begin{pmatrix}0&M(z)\\-M(z)^t&0\end{pmatrix},
 \qquad \det M(z)=0
\]
が $E_L$ の方程式となる。
$E_L$ は滑らかなので $M(z)$ は $E_L$ 上の全ての点で階数二である。
階数一以下なら全ての余因子が零になり、$\det M$ の三つの偏微分が零となるからである。
従って $\mathcal K_L$ と $U_\pm$ の交わりは線束であり、上記の直和分解を与える。
既に得た $\det\mathcal K=\mathcal O_E(-2)$ と、
非自明Galois元が二項を交換することにより、各次数は $-3$ となる。

最後に、体の次数二の拡大により指数が失われるなら、
零cycleを $K$ へ押し下げることから
$\operatorname{ind}(E)\mid2\operatorname{ind}(E_L)$ となる。
一方、平面三次曲線なので $\operatorname{ind}(E_L)\mid3$ である。
$\operatorname{ind}(E)=3$ より $\operatorname{ind}(E_L)=3$ を得る。
周期は指数を割り、周期一の種数一torsorは点を持つので、周期も三である。
$\alpha_L=0$ と一般ファイバーの双有理積表示から残りが従う。
\end{proof}

\subsection{二次分裂体の線形計算とノルムトーラス}

Pl\"ucker関係式を $\Pi$ に制限すれば、分裂体を二変数二次式で求められる。
$\alpha$ が分裂する体上で $\Pi$ の基底を選び、
制限された非零の二次式を $q(a,b)=Aa^2+Bab+Cb^2$ と書く。
全ての制限関係式は $q$ のスカラー倍であり、
\[
 Y_k\simeq\operatorname{Spec}\bigl(k[t]/(t^2-(B^2-4AC))\bigr)
\]
と計算できる。
判別式が平方の場合には、右辺の二次エタール代数は $k\times k$ となる。
ここで式の基礎体は、あくまでこの基底を選べる体 $k$ である。
$\alpha\ne0$ の $K$ 上で $D_\Pi$ を通常の射影直線と同一視してはならない。
基底変更による判別式の変化は行列式の二乗である。

$K$ 上で直接二次代数を書くには、次の別の線形計算が使える。
\[
 \mathfrak a_\Omega=
 \{J\in\operatorname{End}(T):
       \Omega(Ju,v)+\Omega(u,Jv)=0\text{ for all }u,v\}
\]
と置く。$\operatorname{End}(T)$ には慣性が自明に作用するので、
これは $K$ 上の線形方程式である。
この $\operatorname{End}(T)$ は一般には分裂しない次数六の中心単純代数であり、
以下のtraceはその被約traceを意味する。

\begin{proposition}[内在的な二次代数とトーラス]
\label{prop:quadratic-lie-algebra-v28}
$\dim_K\mathfrak a_\Omega=1$ である。
任意の非零元 $J$ に対して
\[
 \operatorname{tr}(J)=0,\qquad J^2=\delta I_6\quad(\delta\in K^\times),
 \qquad L\simeq K[J]\simeq K[t]/(t^2-\delta)
\]
が成立する。
$L^\times$ の $T$ への作用は、$\Omega$ をノルム倍する。
従って $\operatorname{Res}_{L/K}\mathbb G_m/\mathbb G_m$ は
$S\to E$ に忠実に作用する一次元トーラスを与える。
その固定部分は $E_L\to E$ という次数二の有限エタール多重切断である。
\end{proposition}

\begin{proof}
幾何学的な分解 $T=U_+\oplus U_-$ に対して
$\operatorname{diag}(I_3,-I_3)$ は $\mathfrak a_\Omega$ に属する。
従ってその次元は少なくとも一である。
$G=\operatorname{diag}(I_3,I_3)$ における方程式を直接解くと、
\[
 \mathfrak a_\Omega=
 \left\langle\begin{pmatrix}0&I_3\\-I_3&0\end{pmatrix}\right\rangle
\]
である。階数の開性から一般点で次元は一となる。
よって非零元 $J$ は代数閉包上で
$c\operatorname{diag}(I_3,-I_3)$ であり、$J^2=\delta I_6$、
$\operatorname{tr}J=0$ となる。
この恒等式は $K$ に降り、二つの固有空間は $Y$ の二点そのものである。
従って固有値の二次代数 $K[J]$ は $L$ と同一視できる。

$g=rI+sJ$ と書くと、$J$ の定義から
\[
 \Omega(gu,gv)=(r^2-\delta s^2)\Omega(u,v)
       =N_{L/K}(r+s\sqrt\delta)\,\Omega(u,v)
\]
を得る。scalarを除いて射影化すると、上のトーラスが $S$ に作用する。
$L$ 上の各ファイバーでは、二つの核の固有直線を保つ対角トーラスである。
その二つの固定切断はGaloisにより交換されるので、
固定部分全体は $Y\times_KE=E_L$ へ降りる。
\end{proof}

\paragraph{Hesse族での明示式.}
式\eqref{eq:vi-hesse-example-v28}の一変数族では、
$\Pi$ の基底を縮約行列の核から取ると、Pl\"ucker関係式は
\[
 q_t(a,b)=(t-4)a^2+(t-4)ab-b^2=0
\]
となる。従って二次分裂代数は
\[
 \mathbb C(t)\bigl(\sqrt{t(t-4)}\bigr)
\]
である。対応するLie代数の生成元は
\[
 J_t=\begin{pmatrix}
 t-4&0&0&0&2(t-4)&0\\
 0&-t&0&0&0&-2\\
 0&0&t-4&-2&0&0\\
 0&0&-2(t-4)&4-t&0&0\\
 2&0&0&0&4-t&0\\
 0&2t&0&0&0&t
 \end{pmatrix},
 \qquad J_t^2=t(t-4)I_6.
\]
この族では $\det G=t-4$ なので、二次分裂体を単に
$\sqrt{\det G}$ で与えることはできない。
付属の \nolinkurl{isotropic_plucker_v28.py} は、
縮約行列の階数、全二百二十五個のPl\"ucker関係式の制限、
$J_t$ の三つの交代形式に対する恒等式と、その二乗を正確に照合する。
一般モジュライの境界全体に沿うこの二次被覆の分岐は、ここでは決定しない。

\section{逆Gram族の支配性と二つの楕円moduli}
\label{sec:inverse-gram-dominance-v29}

Hermite曲線がnetのみに依存することから、その上の核束もnetだけで決まるとは限らない。
この節では、逆Gram構成が六次元上の一般の交代二形式の三つ組を、基底変換を許して全て実現することを示す。
その結果、Hermite曲線とは別に、そのJacobian上の自由に変化する点が現れる。
以下の支配性は標数零で成立し、幾何学的moduliに関する結論は代数閉体上で述べる。

$\mathcal H=\operatorname{Sym}^2(k^2)\otimes\operatorname{Sym}^2(k^3)$ を
十八次元のブロック対称行列空間とし、
\[
 G=\begin{pmatrix}A&C\\C&B\end{pmatrix},\qquad
 D_i=\operatorname{diag}([e_i]_\times,[e_i]_\times),\qquad
 \Omega_{ij}(G)=D_i^tG^{-1}D_j-D_j^tG^{-1}D_i
\]
と置く。$A,B,C$ は対称な $3\times3$ 行列である。
$\mathcal H^\times=\{\det G\ne0\}$ と書く。

\begin{theorem}[逆Gram構成の全微分階数]
\label{thm:inverse-gram-full-dominance-v29}
射
\begin{align*}
 \mathcal F:\mathcal H^\times\times\operatorname{GL}_6
 &\longrightarrow\bigl(\textstyle\bigwedge^2(k^6)^*\bigr)^{\oplus3},\\
 (G,g)&\longmapsto
 (g^t\Omega_{23}(G)g,g^t\Omega_{31}(G)g,g^t\Omega_{12}(G)g)
\end{align*}
は支配的である。特に、出力の三形式を取り替える $\operatorname{GL}_3$ の作用を
追加せずとも、その四十五次元の標的全体を支配する。
\end{theorem}

\begin{proof}
一点での微分の全射性を有理数体上で検証する。
$g=I_6$ とし、
\[
 A=\begin{pmatrix}0&1&0\\1&0&0\\0&0&1\end{pmatrix},\quad
 C=\begin{pmatrix}1&0&0\\0&0&2\\0&2&0\end{pmatrix},\quad
 B=\begin{pmatrix}0&0&1\\0&1&0\\1&0&0\end{pmatrix}
\]
を取る。このとき $\det G=-3$ である。
微分は
\[
 d(G^{-1})[H]=-G^{-1}HG^{-1},\qquad
 d(g^t\Omega g)[E]=E^t\Omega+\Omega E
\]
により計算する。
四十五行は $(\Omega_{23},\Omega_{31},\Omega_{12})$ の順に、
各交代行列の $i<j$ 成分を辞書式順序で並べる。
五十四列は、$A,C,B$ の順の対称成分 $0\le i\le j\le2$ を辞書式順序で十八本、
続いて $\operatorname{GL}_6$ の $E_{ij}$ を $0\le i,j\le5$ の辞書式順序で三十六本並べる。
この零始まりの列番号のうち
\[
 \begin{gathered}
 0,1,2,3,4,5,6,7,8,9,10,11,12,13,14,15,16,17,\\
 18,19,20,21,22,23,24,25,26,27,28,29,30,31,33,34,\\
 36,37,38,39,40,41,42,43,44,45,46
 \end{gathered}
\]
の四十五列を選ぶと、その小行列式は
\begin{equation}\label{eq:dominance-minor-v29}
 \frac{962948}{1162261467}
 =\frac{2^2\cdot7^2\cdot17^3}{3^{19}}\ne0
\end{equation}
である。

この有限計算は付属の
\nolinkurl{ruled_moduli_dominance_v29.py}で定義から再計算できる。
全ての行を九倍すると最大絶対成分三十六の整数行列となり、
整数の分数なし消去による行列式は
\[
 7231225761666696040510882238572188368556
 =2^2\cdot3^{71}\cdot7^2\cdot17^3
\]
である。従って微分の階数は四十五であり、標数零の微分判定により支配性を得る。
同じ点のPfaffian三次曲線は非零スカラーを除いて
\[
 2z_1^3+z_2^3+z_3^3-7z_1z_2z_3=0
\]
であり、$27\cdot2\cdot1\cdot1-7^3=-289\ne0$ なので滑らかでもある。
\end{proof}

\subsection{行列式線束のmoduli}

一般の交代二形式の三つ組は二つの相補的な被約共通等方三平面を持つ。
それらを順序付けると
\[
 \Omega(z)=\begin{pmatrix}0&M(z)\\-M(z)^t&0\end{pmatrix},\qquad
 0\to\mathcal O_{\mathbb P^2}(-2)^3\xrightarrow{M}
 \mathcal O_{\mathbb P^2}(-1)^3\to i_*\theta\to0
\]
が得られる。$E=\{\det M=0\}$ が滑らかなとき
$\theta\in\operatorname{Pic}^0(E)\setminus\{\mathcal O_E\}$ である。
左右の基底変換は $\theta$ を変えず、二平面の交換は転置により
$\theta\mapsto\theta^{-1}$ を与える。

\begin{corollary}[第二の楕円moduliとnet非因子化]
\label{cor:ruled-two-moduli-v29}
逆Gram族から得る
\[
 (E,\{\theta,\theta^{-1}\})
\]
は二次元の幾何学的moduliを支配する。
従って一般netを固定した五次元の二次曲線パラメータ空間上でも、
$\operatorname{Pic}^0(E)/\{\pm1\}$ における $\theta$ の像は非定数である。
幾何学的一般点で $\theta$ は無限位数であり、複素数体上の非常に一般の点でも同様である。
また、$\theta$ および $\mathbb P_E(\mathcal K)$ の幾何学的同型類はnetのみからは決まらない。
\end{corollary}

\begin{proof}
固定した滑らかな平面三次曲線 $E$ に対し、一次行列式表示の存在と一意性
\cite[Propositions~3.1 and~3.3]{BeauvilleDeterminantal}は、
$3\times3$ の行列 $M$ の左右の基底変換による同値類と
$\operatorname{Pic}^0(E)\setminus\{\mathcal O_E\}$ を対応させる。
従って共通等方三平面を順序付けた一般の交代二形式の三つ組には、曲線 $E$ と点 $\theta$ の
二つの連続moduliがある。順序を忘れる操作は $\theta$ と $\theta^{-1}$ を同一視する有限商である。
次数三の平面埋め込みの線束は、代数閉体上では平行移動で調整できるため、
さらに連続moduliを加えるものではない。
定理\ref{thm:inverse-gram-full-dominance-v29}から、逆Gram族もこの二次元空間を支配する。

Gram--Prym表示の基底変換に対する共変性により、この不変量は
$\mathcal R_3$ の幾何学的同型類に依存する。
一方、$E$ はnetだけから決まり、netのmoduliは一次元である。
もし一般netを固定したとき $\theta$ のKummer像も定数なら、上の対の全像は
高々一次元となり、支配性に反する。

各正整数 $n$ について $\theta^n=\mathcal O_E$ は二次元のmoduli内の真の閉条件である。
従って幾何学的一般点ではどの $n$ に対しても成立せず、
複素数体上では可算個の真閉部分の合併を除いて同じ結論が成立する。
核束は $\mathcal K=\mathcal O_E(-1)\otimes(\theta^{-1}\oplus\theta)$ なので、
その射影化は $\mathbb P_E(\mathcal O_E\oplus\theta^2)$ となる。
二乗写像は有限であるため、この線織面の同型類もnetのみからは決まらない。
\end{proof}

以上は幾何学的moduliに関する主張である。
粗モジュラー基底上に普遍楕円曲線がそのまま存在するとは仮定していない。
また無限位数の結論は、単一の非空Zariski開集合の全ての複素点に対する主張ではない。
有限安定化群の反転と共通等方平面の交換を区別した降下は、次の節で扱う。

\section{符号でねじった二次被覆と楕円曲線への写像}
\label{sec:third-quadratic-cover-v29}

第\ref{sec:intrinsic-quadratic-splitting-v28}節の二次拡大は、共通等方三次元部分空間の一方を選ぶものであった。
ここでは、netの安定化群がJacobianに及ぼす符号作用も残す。
この二つの符号を掛け合わせた第三の二次被覆には、
固定した楕円曲線への非定数有理写像が存在する。
これにより、net写像の有理な五次元一般ファイバーの上に、
幾何学的単有理性を持たない自然な二次被覆が現れる。

\subsection{固定netと二つの符号}

$F=\mathbb C(Y_0(2))$ とし、
\S\ref{sec:modular-fiber-rationality-v28}で構成した $F$ 上のnet代表 $N=N_u$ を固定する。
まず $k=\overline F$ 上で議論する。
二次曲線の空間とその商を
\[
 X=\mathbb P(\operatorname{Sym}^2N_k^*)\simeq\mathbb P_k^5,
 \qquad M=X/G,
 \qquad G\simeq(\mathbb Z/3)^2\rtimes\mathbb Z/2
\]
と置く。以下の有限被覆と有理写像は、必要に応じて稠密開部分へ制限して考える。
$M$ はnet写像 $\Psi$ の幾何学的一般ファイバーと双有理である。
$H\simeq(\mathbb Z/3)^2$ を $G$ の唯一のSylow三部分群、
$\chi:G\to\{\pm1\}$ を商指標とする。

$E_N$ を固定netのHermite曲線、$J=\operatorname{Pic}^0(E_N)$ とする。
共通等方三次元部分空間の一方を選ぶ二次被覆を
\[
 \pi:Y\longrightarrow X
\]
と書き、その交換対合を $\iota$ と書く。
選んだ部分空間を第一成分に置くと、交代行列は
\[
 \Omega(z)=\begin{pmatrix}0&M(z)\\-M(z)^t&0\end{pmatrix}
\]
となる。その余核線束
\[
 \vartheta=\operatorname{coker}
 \bigl(M:\mathcal O_{\mathbb P^2}(-2)^3
                    \longrightarrow\mathcal O_{\mathbb P^2}(-1)^3\bigr)
 \in\operatorname{Pic}^0(E_N)
\]
は基底に依存せず、有理写像 $\vartheta:Y\dashrightarrow J$ を定める。
共通平面の交換は $M$ を転置行列へ移すので
\begin{equation}
 \vartheta(\iota y)=-\vartheta(y)
 \label{eq:theta-deck-sign-v29}
\end{equation}
である。ここでは $J$ の群法を加法で書いた。
この線束と転置の関係は、一次行列式表示の双対性
\cite[Proposition~3.1]{BeauvilleDeterminantal}からも従う。

\begin{lemma}[安定化群の符号と被覆の持上げ]
\label{lem:theta-stabilizer-sign-v29}
$G$ は $Y$ に作用し、その作用は $\iota$ と可換である。
また
\begin{equation}
 \vartheta(gy)=\chi(g)\vartheta(y)
 \label{eq:theta-stabilizer-sign-v29}
\end{equation}
が成立する。
\end{lemma}

\begin{proof}
netの自己同型は交代形式の三つ組とその共通等方部分空間を移す。
従って $Y$ への作用はこの構成の関手性から得られる。
Prymの $\mu_2$ 慣性群は関係空間にスカラーとして作用するため、
Grassmann多様体の点を固定する。
従って、この有限被覆への作用にはスカラーを持ち上げる際の曖昧さが残らない。
二つの共通平面を交換する操作も関手的であり、$G$ の作用と可換である。

$G$ は非退化平面三次曲線 $E_N$ に忠実に作用する。
一般の $j$ 不変量では、種数一曲線の自己同型からJacobianへの作用は
$\{\pm1\}$ に値を取り、その核は平行移動群である。
従って奇数位数の部分群 $H$ は平行移動として作用し、$J$ への作用は自明である。
もし $G$ 全体の $J$ への作用が自明なら、忠実な $G$ の作用は平行移動群に含まれ、
$G$ が非可換であることに反する。
従って $J$ への作用はちょうど $\chi$ である。
余核線束は行列表示の移送と可換なので、\eqref{eq:theta-stabilizer-sign-v29}を得る。
\end{proof}

\begin{lemma}[生の順序付け被覆の既約性]
\label{lem:raw-order-cover-integral-v29}
$Y$ は既約であり、$\vartheta:Y\dashrightarrow J$ は支配的である。
\end{lemma}

\begin{proof}
系\ref{cor:ruled-two-moduli-v29}により、
逆Gram族の像は $(E,\{\vartheta,\vartheta^{-1}\})$ の二次元モジュライを支配する。
一方、$E$ の同型類はnetのみから決まり、netのモジュライは一次元である。
従って、一般netを固定したときのKummer座標
$\{\vartheta,\vartheta^{-1}\}$ は非定数である。

もし $Y\to X$ が分裂していれば、その二つの成分は $X$ と双有理である。
射影空間から楕円曲線への有理写像は定数なので
\cite[I, Proposition~3.10]{MilneAbelianVarieties2008}、
各成分上で $\vartheta$ は定数となる。
二つの値は\eqref{eq:theta-deck-sign-v29}により互いに逆であり、
Kummer座標も定数となって矛盾する。
よって $Y$ は既約である。
同じ非定数性から $\vartheta$ の像は曲線 $J$ で稠密となる。
\end{proof}

\subsection{第三の二次被覆}

$W=Y/H$ と置く。$G/H$ の非自明元が $W$ に誘導する対合を $\rho$ と書く。
これは商群の作用として定まるので $\rho^2=1$ であり、
$G$ の中で特定の持上げを選ぶ必要はない。
$\iota$ は $W$ に降り、$\rho\iota=\iota\rho$ である。

\begin{theorem}[第三の二次被覆の非単有理性]
\label{thm:third-quadratic-cover-v29}
$W\to M$ は一般点上で群 $(\mathbb Z/2)^2$ のGalois被覆である。
その三つの中間二次被覆は
\[
 M_\chi=X/H=W/\langle\iota\rangle,\qquad
 M_L=Y/G=W/\langle\rho\rangle,\qquad
 Z=W/\langle\rho\iota\rangle
\]
である。$M_L$ は共通平面から得られる二次拡大 $L/K$ の
幾何学的一般netファイバーへの引戻しである。
第三の被覆 $Z$ には支配的有理写像
\[
 \vartheta_Z:Z\dashrightarrow J
\]
が存在する。
従って $Z$ の任意の滑らかな射影モデル $\overline Z$ は
\[
 q(\overline Z)=h^1(\overline Z,\mathcal O_{\overline Z})\ge1
\]
を満たす。特に $Z$ は単有理でなく、安定有理でもない。
同じ結論は $Y$ と $W$ の滑らかな射影モデルにも成立する。
\end{theorem}

\begin{proof}
$Y$ の既約性と $G$ の $X$ への忠実な作用により、
$G\times\langle\iota\rangle$ は $Y$ に忠実に作用する。
一般点では作用は自由であり、$H$ で割ると
$(G/H)\times\langle\iota\rangle\simeq(\mathbb Z/2)^2$ を得る。
これが三つの中間被覆とその次数を与える。

$H$ は $J$ に自明に作用するため、$\vartheta$ は $W$ に降りる。
\eqref{eq:theta-deck-sign-v29}と\eqref{eq:theta-stabilizer-sign-v29}によれば
$\iota$ と $\rho$ はともに $\vartheta$ の符号を反転する。
従って積 $\rho\iota$ は $\vartheta$ を保ち、$Z$ への降下を得る。
元の写像が支配的なので $\vartheta_Z$ も支配的である。

滑らかな多様体からAbel多様体への有理写像は全体に延長する
\cite[I, Theorem~3.2]{MilneAbelianVarieties2008}。
従って $\overline Z\to J$ は非定数の射であり、
Albanese多様体の普遍性から $\operatorname{Alb}(\overline Z)\to J$ は全射となる。
標数零では $\dim\operatorname{Alb}(\overline Z)=q(\overline Z)$ なので $q\ge1$ である。
単有理多様体からAbel多様体への有理写像は定数であるため、$Z$ は単有理でない。
安定有理性は単有理性を含むので、後者も排除される。
$Y$ と $W$ についても、同じ支配的写像を用いればよい。
\end{proof}

\begin{corollary}[二つの二次指標の独立性]
\label{cor:quadratic-character-independence-v29}
$K_0=k(M)$ とし、$k(M_L)=K_0(\sqrt\delta)$、
$k(M_\chi)=K_0(\sqrt\epsilon)$ と書く。
この二つの二次拡大は相異なり、
\[
 k(W)=K_0(\sqrt\delta,\sqrt\epsilon),\qquad
 k(Z)=K_0(\sqrt{\delta\epsilon}).
\]
また $\vartheta_Z$ が定める $J(k(Z))$ の点は無限位数を持つ。
\end{corollary}

\begin{proof}
二つの拡大が一致すれば、$K_0(\sqrt\epsilon)\subset k(X)$ なので
$Y\to X$ は分裂し、補題\ref{lem:raw-order-cover-integral-v29}に反する。
四次の合成体の第三中間体は表示通りである。
さらに $n\vartheta_Z=0$ なら、$\vartheta_Z$ の像は有限群 $J[n]$ に含まれ、
支配性に反する。
\end{proof}

\subsection{Kummer被覆としての解釈と基礎体への降下}

二つの符号で商を取ると、$\vartheta$ のKummer座標は $M$ に降りる。
従って非定数有理写像
\[
 \kappa:M\dashrightarrow J/\{\pm1\}\simeq\mathbb P^1
\]
が存在する。上の構成は
\begin{equation}
 Z\dashsim \operatorname{Norm}
 \bigl(M\times_{J/\{\pm1\}}J\bigr)
 \label{eq:third-cover-kummer-pullback-v29}
\end{equation}
を与える。ここで右辺は関数体の意味での正規化である。
実際、$W$ の関数体内で $k(J)$ の元を固定するのは
$\{1,\rho\iota\}$ であり、二次被覆の両辺の関数体が一致する。

この構成は、固定したnet代表 $N_u/F$ に対して $F$ 上へ降りる。
$H$ は $G$ の特徴的部分群スキームであり、$G/H$ は位数二の定数群スキームである。
共通平面の有限被覆 $Y$ とその交換 $\iota$ も $F$ 上で定まる。
線束 $\vartheta$ の構成をPrymのgerbe上で行う場合、その余核線束 $\Theta$ は重み一である。
また、選んだ階数三の重み一部分束 $U_+$ の行列式 $\ell=\det U_+$ も重み一である。
従って $\Theta\otimes\operatorname{pr}_Y^*\ell^{-1}$ は重み零となり、
適切な稠密開部分上の $E_N\times Y$ に降下する。
基底からの線束のテンソル積は相対Picard classを変えないため、
これにより $F$ 上の有理写像 $\vartheta:Y\dashrightarrow J$ が定義される。
任意の基底上でgerbe全体が中立であるとは仮定しない。
その同変性は代数閉包上の等式から降りるため、
$W,\rho,Z,\vartheta_Z,\kappa$ および\eqref{eq:third-cover-kummer-pullback-v29}も
$F$ 上で得られる。

特に、$K=F(M)=\mathbb C(\mathcal R_3)$ に対して
\[
 K_\chi=F(X)^H,
 \qquad [LK_\chi:K]=4,
 \qquad F(Z)=K(\sqrt{\delta\epsilon})
\]
である。一般Hermite曲線のJacobianは、固定した $J_F$ を
$G$ の符号作用 $\chi$ で降ろした二次ねじれである。
このねじれを忘れて $L$ だけを取ると、残る安定化群の対合が
$\vartheta$ を反転する。第三の二次被覆は、共通平面の交換とこの反転を相殺する。

\subsection{符号被覆の判別式}

\S\ref{sec:modular-fiber-rationality-v28}の表現
$V'=V\otimes\chi$ を用いる。
群が分裂する体上のHesse座標では
\[
 p=x^3+y^3+z^3,\qquad q=xyz,\qquad
 r=x^3y^3+y^3z^3+z^3x^3
\]
が不変式を生成する。
一方
\[
 d=(x^3-y^3)(y^3-z^3)(z^3-x^3)
\]
は符号指標 $\chi$ に従って変換する。
三変数 $x^3,y^3,z^3$ のVandermonde式から
\begin{equation}
 d^2=p^2r^2-4r^3-4p^3q^3-27q^6+18prq^3
 \label{eq:sign-cover-discriminant-homogeneous-v29}
\end{equation}
を得る。
階数三ベクトル束の商の表示において、$G/H$ の被覆はその底からの引戻しである。
従って $q\ne0$ 上で
\[
 a=p/q,\qquad b=r/q^2,\qquad w=d/q^3
\]
と置くと、符号被覆の関数体は
\begin{equation}
 w^2=\Delta(a,b),\qquad
 \Delta(a,b)=a^2b^2-4b^3-4a^3-27+18ab
 \label{eq:sign-cover-discriminant-affine-v29}
\end{equation}
で与えられる。この座標では、第三の被覆は
\[
 K_0\bigl(\sqrt{\delta\Delta(a,b)}\bigr)
\]
である。

$F$ 上でこの特定のHesse座標を選べるとは限らないが、符号被覆自体は降りる。
実際、$\chi$ に従って変換する多項式全体は、代数閉包上で
$d\,\overline F[p,q,r]$ である。
従って次数九の半不変式空間は $F$ 上一次元であり、非零元 $d_0$ を選べる。
不変式の $F$ 上の生成元 $p_0,q_0,r_0$ に対して
$d_0^2=D_0(p_0,q_0,r_0)$ という重み付き次数十八の関係を得る。
その平方類 $D_0/q_0^6$ が $\epsilon$ を表し、
\eqref{eq:sign-cover-discriminant-affine-v29}はその分裂座標での表示である。

\paragraph{次節への接続と結論の範囲.}
この節では幾何学的一般netファイバーについて $q\ge1$ を得た。
次節ではKummer写像を三次pencilの関数で明示し、$q=1$ とAlbanese曲線を決定する。
このAlbanese曲線から $J$ への有限部分は二次同種写像であり、
$\vartheta_Z$ の幾何学的一般ファイバーは二つの連結成分を持つ。
従って $J$ をそのまま連結ファイバーの底とみなしてはならない。
全六次元空間の不正則数や非有理性、および元の共通平面被覆 $M_L$ の有理性は別の問題である。

\paragraph{再現入口.}
\nolinkurl{modular_quadratic_covers_v29.py}は、単項行列群の位数と符号部分群、
半不変式 $d$ の変換則、\eqref{eq:sign-cover-discriminant-homogeneous-v29}、
および\eqref{eq:sign-cover-discriminant-affine-v29}を厳密演算で検証する。
被覆の既約性、群スキーム上の降下、楕円曲線への支配性と非単有理性は、
この節の証明と系\ref{cor:ruled-two-moduli-v29}に依存する。

\section{三次pencilによるKummer写像と第三被覆のAlbanese}
\label{sec:kummer-albanese-v30}

本節では幾何学的一般netを固定する。基礎体 $k$ は標数零の代数閉体とし、
相互に双対な二次曲線 $D\subset\mathbb P(N)$ と $D^\vee\subset\mathbb P(N^*)$ のうち、
以下では $D^\vee$ の係数を座標として $X=\mathbb P^5$ と書く。
前節で $D$ の係数を用いた $X$ とは双対化による双有理同値で同定する。
一般netの自己同型群を
$G=H\rtimes\langle\rho\rangle$、$H\simeq(\mathbb Z/3)^2$ と書く。
二共通等方平面の順序付けを $Y\to X$、その対合を $\iota$ とし、
\[
 W=Y/H,\qquad Z=W/\langle\rho\iota\rangle
\]
を用いる。以下の有理写像の射影モデルは、必要に応じて不定性を解消して取る。
ここでのAlbaneseと不正則数の主張は、固定された幾何学的一般netの五次元空間に関するものである。
全六次元モジュライ上の被覆に同じ不正則数を主張するものではない。

\subsection{三次の有理関数と明示的な二次同種写像}

netの有限なパラメータ被覆として
\begin{equation}\label{eq:net-m-v30}
 A_0=\begin{pmatrix}0&1&0\\1&0&0\\0&0&1\end{pmatrix},\quad
 A_1=\begin{pmatrix}1&0&0\\0&0&2\\0&2&0\end{pmatrix},\quad
 A_2(m)=\begin{pmatrix}0&0&1\\0&m&0\\1&0&0\end{pmatrix}
\end{equation}
を用いる。$m$ は一般とする。この族のnet判別曲線とHermite曲線はそれぞれ
\begin{align}
 \Delta_m &: -\lambda_0^3-4\lambda_1^3-m\lambda_2^3
                  +(m+4)\lambda_0\lambda_1\lambda_2=0,\label{eq:delta-m-v30}\\
 E_m &:2mz_0^3+z_1^3+mz_2^3+(m-8)z_0z_1z_2=0.\label{eq:hermite-m-v30}
\end{align}
いずれも $m(m-2)(m+16)\ne0$ なら滑らかである。
この一変数族はnetの一次元モジュライを支配するため、幾何学的一般netの議論には十分である。

順序付き基底 $(a,b,c)$ の $\langle A_0,A_1,A_2(m)\rangle$ に関する係数行列を
$L$ とし、係数列ベクトルを $a_*,b_*,c_*$ と書く。次を定める。
\begin{equation}\label{eq:cubic-R-v30}
 C=a_*c_*^t+c_*a_*^t-2b_*b_*^t,\qquad
 I_m(C)=\det\begin{pmatrix}a&b\\b&c\end{pmatrix},\qquad
 R=\frac{\det\Gamma}{(\det L)^2}=\frac{2I_m(C)}{\det C}.
\end{equation}
ここで $a,b,c$ は $3\times3$ 対称行列であり、
$\Gamma=\left(\begin{smallmatrix}a&b\\b&c\end{smallmatrix}\right)$ と書く。
$C$ は双対二次曲線 $ac-b^2=0$ の係数行列の二倍である。
$I_m$ は $C$ の六係数の三次多項式であり、従って $R$ は $X$ 上の三次pencilを定める。
式\eqref{eq:cubic-R-v30}は二次元基底の変更に不変である。

\begin{theorem}[Kummer写像、第三被覆とAlbanese]
\label{thm:kummer-albanese-v30}
$q=(m-2)^3$、$a_m=m^2-40m-32$ と置く。
Hermite曲線のJacobianのWeierstrassモデルとして
\begin{equation}\label{eq:Jm-v30}
 J_m:\quad y_J^2=x_J^3-2a_mx_J^2+m(m+16)^3x_J
\end{equation}
を取れる。この座標で行列式線束のKummer写像は
\begin{equation}\label{eq:kappa-R-v30}
 \boxed{\quad\kappa(R)=-4R+a_m+\frac{8q}{R}\quad}
\end{equation}
である。次の楕円曲線を定める。
\begin{equation}\label{eq:Bm-v30}
 B_m:\quad v^2=F_m(R),\qquad
 F_m(R)=-R(4R^2-a_mR-8q).
\end{equation}
このとき以下が成立する。
\begin{enumerate}
\item 共通平面順序付けの関数体は
\[
 k(Y)=k(X)\bigl(\sqrt{F_m(R)}\bigr)
\]
である。従って $Y$ は $X\times_{\mathbb P^1_R}B_m$ の正規化と双有理である。
\item $Y,W,Z$ はすべて $B_m$ への支配的有理写像を持ち、その幾何学的一般ファイバーは
四次元、幾何学的整かつ単有理である。従って各幾何学的一般ファイバーの
滑らかな射影モデルは有理連結である。
\item 滑らかな射影モデルについて
\begin{equation}\label{eq:Alb-v30}
 \operatorname{Alb}(Y)\simeq\operatorname{Alb}(W)
 \simeq\operatorname{Alb}(Z)\simeq B_m,
 \qquad q(Y)=q(W)=q(Z)=1.
\end{equation}
幾何学的に $B_m\simeq\operatorname{Jac}(\Delta_m)$ である。
\item $\theta_Z:Z\dashrightarrow J_m$ のStein分解の有限部分は二次同種写像
\begin{equation}\label{eq:phi-m-v30}
 \begin{split}
 \varphi:B_m&\longrightarrow J_m,\\
 (R,v)&\longmapsto\left(
 -4R+a_m+\frac{8q}{R},\,
 -\frac{4v(R^2+2q)}{R^2}\right).
 \end{split}
\end{equation}
したがって $\theta_Z$ の幾何学的一般ファイバーは、ちょうど二つの連結成分を持つ。
各成分は四次元で単有理である。
\end{enumerate}
\end{theorem}

\paragraph{四つの三次超曲面を分岐因子とするモデル.}
$D=\det C$、$I=I_m(C)$ と略記する。式\eqref{eq:Bm-v30}に $R=2I/D$ を代入し、
重み六の座標 $w=vD^2$ を用いると
\begin{equation}\label{eq:four-cubics-double-cover-v30}
 w^2=-4ID(8I^2-a_mID-4qD^2)
\end{equation}
を得る。この式は $Y$ の正規な有限二重被覆モデルを与える。
平方根の下の式は三次pencilの四つの異なるmemberの積であり、総次数は十二である。
一般の $m$ でこの積は平方因子を持たない。これは $m=1$ で全偏微分との最大公約数が一である
exactな証明書と、固定次数の多項式の平方因子を持たない条件の開性から従う。
従ってこの四memberの和が分岐因子である。
被覆は超曲面なので $S_2$ を満たし、被約な分岐因子により余次元一で正則だから正規である。
これらは共通の基底部分集合を持つので、一般の滑らかな十二次分岐因子の場合の幾何学を
そのまま当てはめることはできない。

この定理は $Z\dashsim B_m\times\mathbb P^4$ や、四次元ファイバーの有理性を主張しない。
また $k(J_m)$ 上の一般ファイバーは整であり、その定数体が二次拡大 $k(B_m)$ である。
二成分となるのは幾何学的一般ファイバーである。

\subsection{三次不変式と交換子のtrace}

まず $I_m$ の存在と規約を確認する。
二次元基底の変更 $U\in\operatorname{GL}_2$ により $\Gamma$ は $U\otimes I_3$ で合同変換され、
$\det\Gamma$ は $(\det U)^6$ 倍になる。
同時に $\det L$ は $(\det U)^3$ 倍になるから、$R$ は基底に依存しない。
$\operatorname{SL}_2\to\operatorname{SO}_3$ と三次元直交群の不変式を用いると、
次数六の $\det\Gamma$ はGram係数 $C$ の三次式となる。
第\ref{subsec:explicit-cubic-v30}節の全係数によっても、この点は独立に検証できる。
また
\[
 \det C=2(\det L)^2.
\]
階数一の行列 $C=ww^t$ に対しては
\begin{equation}\label{eq:rankone-I-v30}
 I_m(ww^t)=\frac18\det\!\left(\sum w_iA_i\right)^2.
\end{equation}
実際、$a=c=(\sum w_iA_i)/\sqrt2$、$b=0$ と置けばよい。
従って $I_m$ は $\det C$ の定数倍ではなく、$R$ は非定数である。

行列はPrym gerbe上で枠を選んで計算する。
以下の $h$ は $\operatorname{End}(T)$ の被約traceであり、基底不変なので粗空間へ降りる。
三つの交代形式を $\Omega_0=\Omega_{23}$、$\Omega_1=\Omega_{31}$、
$\Omega_2=\Omega_{12}$ と書き、$\Omega_0$ が可逆な開部分で
\begin{equation}\label{eq:commutator-v30}
 A=\Omega_0^{-1}\Omega_1,\quad B=\Omega_0^{-1}\Omega_2,\quad
 K=[A,B],\quad h=\operatorname{tr}(K^2),\quad
 \mathscr J=K^3-\frac h4K
\end{equation}
を定める。逆Gram構成に対する主要な恒等式は
\begin{equation}\label{eq:trace-R-v30}
 \boxed{\quad h=
 \frac{-3R+(m-20)(m-2)+6(m-2)^3/R}{m^2}.\quad}
\end{equation}

式\eqref{eq:trace-R-v30}は六独立変数上の多項式恒等式として検証する。
その際に使う稠密sliceを明示しておく。
\[
 P=A_0+\alpha A_1+\beta A_2,\quad
 Q=\gamma A_1+\delta A_2,\quad T=\varepsilon A_2,
 \qquad (a,b,c)=(P+Q,T,P-Q).
\]
これは二次形式 $P^2-Q^2-T^2$ の三角分解であり、$\operatorname{SO}_3$ の作用と
共通スカラーを加えれば全基底空間の稠密開部分を覆う。
両辺はこの作用に不変であるから、このsliceでの恒等式は一般の場合を証明する。
具体的な証明書では、$\Gamma$ の余因子行列から $\Omega_i$ の分子を作り、
$\Omega_0^{-1}$ を四次Pfaffianで書き、分母を払って全係数を比較する。
変数は $\alpha,\beta,\gamma,\delta,\varepsilon,m$ の六個であり、特殊点での数値確認ではない。

\subsection{Moore標準形からのKummer座標}

共通平面を順序付けると
\[
 \Omega(z)=\begin{pmatrix}0&M(z)\\-M(z)^t&0\end{pmatrix}
\]
となる。$C_3=[M_0^{-1}M_1,M_0^{-1}M_2]$ と置くと、$K$ の二ブロックは
$C_3$ と $-C_3^t$ にそれぞれ相似である。
$\operatorname{tr}C_3=0$ であり、Cayley--Hamilton恒等式より
\[
 C_3^3-\tfrac12\operatorname{tr}(C_3^2)C_3=\det(C_3)I_3.
\]
従って $\mathscr J$ は二共通平面上で互いに反対のスカラーとして作用し、
その固有値を $\ell$ とすれば $\mathscr J^2=\ell^2I_6$ となる。
一般点で $\ell\ne0$ だから、これが順序付けの二次拡大を生成する。

このスカラーと行列式線束との関係はMoore標準形で求められる。
Hesse三次曲線 $z_0^3+z_1^3+z_2^3-pz_0z_1z_2=0$ に対し、
Buchweitz--Pavlovの標準形\cite[Theorem A and Remark 1.4(b)]{BuchweitzPavlovMoore}は
\[
 M_u(z)=\begin{pmatrix}
 u_0z_0&u_1z_2&u_2z_1\\
 u_1z_1&u_2z_0&u_0z_2\\
 u_2z_2&u_0z_1&u_1z_0
 \end{pmatrix},\qquad
 p=\frac{u_0^3+u_1^3+u_2^3}{u_0u_1u_2}
\]
を与える。全ての一般の一次行列式表示がこの形に同値である。
このとき $C_3$ は対角行列で、その固有値は
\[
 \frac{u_2^3-u_1^3}{u_0u_1u_2},\quad
 \frac{u_1^3-u_0^3}{u_0u_1u_2},\quad
 \frac{u_0^3-u_2^3}{u_0u_1u_2}
\]
である。
\[
 r=\frac{u_0^3u_1^3+u_1^3u_2^3+u_2^3u_0^3}{(u_0u_1u_2)^2},\qquad d=\det C_3
\]
と置けば
\begin{equation}\label{eq:Moore-trace-v30}
 h=4(p^2-3r),\qquad
 d^2=p^2r^2-4r^3-4p^3-27+18pr.
\end{equation}
第二式は三つの根 $u_i^3/(u_0u_1u_2)$ を持つ三次式の判別式である。

ここで得た楕円曲線は、行列式線束のJacobianそのものである。
標準形の同値関係は $u-u'\in E[3]$ であり、余核線束への写像は三倍写像に一致する。
一方、$(r,d)$ の関数体は同じ $E[3]$ の不変体である。
実際、一般の $(r,d)$ から三次式の三つの根を復元する偶置換の選択は三通りであり、
その立方根を射影座標として復元する選択も三通りである。
この九個の点は、対角位数三作用と巡回置換が生成する $E[3]$ の軌道をなす。
また $u_0u_1u_2=0$ の九変曲点は自明線束へ行き、$(r,d)$ では無限遠点へ行く。
したがってこの同定は原点を保つ。
転置は $d$ の符号を変え、線束の双対化に対応する。

式\eqref{eq:hermite-m-v30}をHesse形へ対角変換して
\eqref{eq:Moore-trace-v30}を戻すと
\begin{equation}\label{eq:commutator-Weierstrass-v30}
 x_J=\frac{4m^2h-(m+16)^2}{3},\qquad y_J=32m^3\ell,
 \qquad
 y_J^2=x_J(x_J^2-2a_mx_J+m(m+16)^3)
\end{equation}
となる。この同定には原点や行列式線束の有限倍の不確定性は残らない。
残る符号は二共通平面の順序に対応する $y_J$ の符号だけである。
式\eqref{eq:trace-R-v30}を代入すると\eqref{eq:kappa-R-v30}を得る。
さらに
\begin{equation}\label{eq:lie-square-factor-v30}
 \ell^2=F_m(R)\frac{(R^2+2q)^2}{64m^6R^4}.
\end{equation}
従って $v=-8m^3R^2\ell/(R^2+2q)$ と置けば\eqref{eq:Bm-v30}が得られ、
\eqref{eq:phi-m-v30}も従う。

$B_m$ で $x=-4R,y=4v$ と置くと
\[
 y^2=x^3+a_mx^2-32qx.
\]
式\eqref{eq:phi-m-v30}はこの曲線の $(0,0)$ に関する二次同種写像であり、
\[
 a_m^2+128q=m(m+16)^3
\]
がその行先の式を与える。Hesse三次式の不変量との比較により
$B_m$ と $\Delta_m$ のJacobianの $j$ 不変量は一致する。
$k$ が代数閉体なので $B_m\simeq\operatorname{Jac}(\Delta_m)$ となる。
この $j$ の比較だけをもって、元のモジュライ関数体上の捻りまで同定するものではない。

\subsection{有限群の作用と幾何学的一般ファイバー}

$R$ が $G$ に不変であることを確認する。
netをHesse標準形に移すと、その一般の自己同型群は巡回置換、
$\operatorname{diag}(1,\omega,\omega^2)$、二座標の交換で生成される。
$\operatorname{SL}_3$ への持上げに対し、netの三次元基底表現の行列式は
前二者で $1$、最後の交換で $-1$ である。
従ってその行列式は符号指標 $\chi$ であり、平方は常に $1$ である。
$\det\Gamma$ はこの持上げに不変で、$(\det L)^2$ も不変となるから、$R$ は $G$ に不変である。

$B_m\to\mathbb P^1_R$ は二次被覆なので、$G$ の各元の作用は恒等または反転である。
$\theta(gy)=\chi(g)\theta(y)$ と $\theta=\varphi\circ(Y\to B_m)$ から、
$g$ は $B_m$ 上で正確に $\chi(g)$ として作用する。
同様に $\iota$ は $B_m$ の反転である。
従って $H$ は $B_m$ を固定し、$\rho\iota$ も $B_m$ を固定する。
これにより $Y,W,Z$ から $B_m$ への写像を得る。

$R=r$ の閉包は $\mathbb P^5$ の三次超曲面
\begin{equation}\label{eq:cubic-fiber-v30}
 X_r:\quad 2I_m(C)-r\det C=0
\end{equation}
である。このpencilの無限遠のmember $\det C=0$ は、階数二以下の対称三次行列の
既約な行列式超曲面である。またその六偏微分は線形独立である。
固定次数の可約多項式の集合は閉集合であり、偏微分の線形独立性は開条件なので、
一般の $X_r$ は幾何学的整で、錐ではない。
代数閉体上の四次元三次超曲面であるから、Koll\'arの単有理性定理
\cite[Theorem 2]{KollarCubicUnirational}を適用できる。
従って $X_r$ は幾何学的に単有理である。
この論証は $X_r$ の滑らかさを仮定しない。

$Y\to B_m$ の幾何学的一般ファイバーはこの $X_r$ と双有理である。
$W,Z$ のファイバーはその有限商であるから、引き続き幾何学的整かつ単有理である。
標数零では滑らかな射影単有理多様体は有理連結である。
特に $Y,W,Z\to B_m$ は幾何学的に連結な一般ファイバーを持つ。
$\varphi$ は次数二なので、これが $\theta_Z$ のStein分解を与える。

\subsection{Albaneseの同定}

$T$ を $Y,W,Z$ のいずれかの滑らかな射影モデルとし、
$\beta:T\to B_m$ を上の有理写像の解消とする。
幾何学的一般ファイバーは単有理であるから、そのAbel多様体への任意の射は定数である。
従ってAlbanese射 $a_T:T\to\operatorname{Alb}(T)$ は $\beta$ を通じて因数分解する。
一般ファイバーの幾何学的連結性と標数零から、これは基底曲線上の射
$g:B_m\to\operatorname{Alb}(T)$ を与える。
原点を選び、$\beta$ が誘導する準同型を $u:\operatorname{Alb}(T)\to B_m$ と書くと
\[
 u\circ g=\operatorname{id}_{B_m}.
\]
原点を保つ楕円曲線からAbel多様体への射 $g$ は準同型である。
さらに $g(B_m)$ は $a_T(T)$ を含み、Albaneseの生成性から $g$ は全射である。
従って $u,g$ は互いに逆な同型である。これが\eqref{eq:Alb-v30}を証明する。


\begin{corollary}[三つの二次被覆の不正則数の区別]
\label{cor:three-cover-q-v30}
$M=X/G$、$M_\chi=X/H$、$M_L=Y/G$ を前節の三つの二次被覆の記号とする。
固定された幾何学的一般net上で、滑らかな射影モデルの不正則数は
\[
 q(M)=q(M_\chi)=q(M_L)=0,\qquad q(Z)=1
\]
である。この結論は $M_L$ の有理性を判定しない。
\end{corollary}

\begin{proof}
$M$ と $M_\chi$ は $\mathbb P^5$ の有限商であるから単有理であり、不正則数は零である。
$M_L$ からAbel多様体への射を $Y$ へ引き戻すと、$\operatorname{Alb}(Y)=B_m$ を通じて因数分解する。
この射は $G/H$ に不変であり、$G/H$ は $B_m$ 上で反転として作用する。
従って基点を合わせた準同型 $f$ は $f=-f$ を満たす。
連結な楕円曲線からAbel多様体への準同型の像は連結なので、$2f=0$ なら $f=0$ である。
これが $q(M_L)=0$ を与える。
\end{proof}

\paragraph{再現と残る範囲.}
\nolinkurl{kummer_cubic_pencil_v30.py} は
\eqref{eq:trace-R-v30}の六変数多項式恒等式、Moore交換子、
\eqref{eq:commutator-Weierstrass-v30}、\eqref{eq:lie-square-factor-v30}、
二次同種写像、および二つの $j$ の照合をexactに再検証する。
三次式 $I_m$ 自体は十変数 $(L,m)$ の全展開による別証明書を持つ。
四つの分岐memberの被約性も、全偏微分との最大公約数を用いて検証する。
幾何学的一般ファイバーの単有理性とAlbaneseの同定は上の構造的証明による。
残る問題は、有限商である四次元ファイバーの有理性、
このAlbanese fibrationの双有理自明性、および元のモジュライ関数体上での算術的捻りの同定である。

\subsection{三次式の全係数}
\label{subsec:explicit-cubic-v30}
$C=(c_{ij})_{0\le i,j\le2}$ を対称行列とする。
式\eqref{eq:cubic-R-v30}の規約では
\[
 8I_m(C)=I_0(C)+mI_1(C)+m^2I_2(C)
\]
であり、各係数は次の通りである。
\begin{align*}
 I_0={}&c_{00}^3+8c_{00}^2c_{12}-16c_{00}c_{01}c_{02}-24c_{00}c_{01}c_{11}\\
 &+16c_{00}c_{12}^2+32c_{01}^3+16c_{01}^2c_{22}-32c_{01}c_{02}c_{12}\\
 &-64c_{01}c_{11}c_{12}+16c_{02}^2c_{11}+32c_{02}c_{11}^2+16c_{11}^3,\\[.4em]
 I_1={}&2c_{00}^2c_{12}-4c_{00}c_{01}c_{02}-6c_{00}c_{02}c_{22}-8c_{00}c_{11}c_{22}\\
 &+16c_{00}c_{12}^2+16c_{01}^2c_{22}-32c_{01}c_{02}c_{12}-16c_{01}c_{11}c_{12}\\
 &+8c_{01}c_{22}^2+8c_{02}^3+16c_{02}^2c_{11}+8c_{02}c_{11}^2\\
 &-16c_{02}c_{12}c_{22}-24c_{11}c_{12}c_{22}+32c_{12}^3,\\[.4em]
 I_2={}&c_{00}c_{11}c_{22}+2c_{01}c_{22}^2-4c_{02}c_{12}c_{22}+c_{22}^3.
\end{align*}
\nolinkurl{modular_cubic_pencil_v30.py} は、$L$ の九係数と $m$ を独立変数として
$I_m(a_*c_*^t+c_*a_*^t-2b_*b_*^t)=\det\Gamma$ を全係数で照合する。
したがってこの式は補間による予想ではなく、多項式恒等式である。


\section{検証と適用範囲}

本稿の主張は、一般点における同型・双有理同型と、
それらを保つ非空開部分に関するものである。
線束の無限位数性については幾何学的一般点、または非常に一般の複素点という範囲を用いる。
特に、net写像の相対有理性と、自然な二次曲線の非分裂性は両立する。
安定化群の商を取ることでモジュライの関数体が有理になっても、
そこに付随する射影表現の持上げや種数一曲線の原点が自動的に得られるわけではない。

有限計算は、一般性を保証する階数と、証明で用いた多項式恒等式を確認する。
降下、支配性、指数の下界、関数体の純超越性を数値標本から推測するものではない。
付属コードの通常実行は保存した証明書との読取り専用照合である。
\begin{longtable}{@{}p{.45\linewidth}p{.50\linewidth}@{}}
\toprule 検証コード&確認する内容\\\midrule
\nolinkurl{modular_fiber_rationality_v28.py}&位数五十四の単項群、位数十八の射影像、二次式表現、net正規形と $j$ 不変量\\
\nolinkurl{hermite_index_three_v27.py}&三次捩れ族の二つの曲線、滑らかさの因子、支配性を与える十八次小行列式\\
\nolinkurl{brauer_conic_cocycle_v27.py}&Prym表示の半線形降下と四元数classのcocycle恒等式\\
\nolinkurl{isotropic_plucker_v28.py}&縮約の階数、全Pl\"ucker二次関係式の制限、二次代数とノルム恒等式\\
\nolinkurl{ruled_moduli_dominance_v29.py}&逆Gram構成の四十五次小行列式と交代二形式の全ての三つ組への支配性の有限入力\\
\nolinkurl{modular_quadratic_covers_v29.py}&符号部分群、Vandermonde判別式、二つの二次指標の作用表\\
\nolinkurl{modular_cubic_pencil_v30.py}&三次式 $I_m$ の十変数全係数恒等式と行列式memberの非錐性\\
\nolinkurl{kummer_cubic_pencil_v30.py}&六変数trace恒等式、Moore標準形、二次同種写像、$j$ 不変量、四member分岐因子の被約性\\
\bottomrule
\end{longtable}

第三の被覆のAlbanese曲線上の幾何学的一般ファイバーは単有理であるが、
その有理性やAlbanese写像の双有理自明性は未決である。
元の共通平面被覆 $L/K$ 自体の有理性、およびHesse被覆から元のモジュラー関数体への
$B_m$ の算術的捻りの同定も未決である。
幾何学的一般netファイバーの非単有理性を、全六次元空間の非有理性へ移すことはしない。
さらに、二次拡大 $L/K$ のモジュライ境界に沿う分岐、
ねじれを保持した相対コンパクト化、および標識を忘れた関係代数のモジュライへの降下も未解決である。
本稿の三次曲線と二次曲線は一般関数体上で定義されるが、
任意の境界モデルに対して同じ形の束が延長されることまでは主張しない。

\clearpage
\begin{thebibliography}{99}
\bibitem{PaperV}
\emph{四次Prymデータに付随する線織面と種数二十K3の境界周期},
本研究の別稿、第30版、2026年。
\nolinkurl{self_opposite_prym_fibrations_ja_v30.pdf}.

\bibitem{DolgachevRationality2008}
I.~V.~Dolgachev, \emph{Rationality of $\mathcal R_2$ and $\mathcal R_3$},
Pure Appl. Math. Q.\ \textbf{4} (2008), no.~2, 501--508.
\url{https://doi.org/10.4310/PAMQ.2008.v4.n2.a8}.

\bibitem{BeauvilleDeterminantal}
A.~Beauville, \emph{Determinantal hypersurfaces},
Michigan Math. J.\ \textbf{48} (2000), 39--64.
\url{https://arxiv.org/abs/math/9910030}.

\bibitem{DolgachevCAG}
I.~V.~Dolgachev, \emph{Classical Algebraic Geometry: A Modern View},
Cambridge University Press, 2012.
著者公開原稿、\S3.2.2 and Exercise~3.12.
\url{https://sites.lsa.umich.edu/idolga/wp-content/uploads/sites/1334/2024/08/CAG.21.pdf}.
\url{https://doi.org/10.1017/CBO9781139084437}.

\bibitem{DomokosFeherRimanyi2013}
M.~Domokos, L.~M.~Feh\'er and R.~Rim\'anyi,
\emph{Equivariant and invariant theory of nets of conics with an application to Thom polynomials},
J. Singul.\ \textbf{7} (2013), 1--20.
\url{https://doi.org/10.5427/jsing.2013.7a}.

\bibitem{ArtebaniDolgachevHesse}
M.~Artebani and I.~Dolgachev, \emph{The Hesse pencil of plane cubic curves},
Enseign. Math.\ \textbf{55} (2009), no.~3/4, 235--273.
\url{https://doi.org/10.4171/LEM/55-3-3}.

\bibitem{BruinPrymGenusThree}
N.~Bruin, \emph{The arithmetic of Prym varieties in genus 3},
Compos. Math.\ \textbf{144} (2008), 317--338.
\url{https://doi.org/10.1112/S0010437X07003314}.
\url{https://arxiv.org/abs/math/0408069}.

\bibitem{StacksBrauerV27}
The Stacks Project Authors, \emph{The Stacks Project}, Tags 03R1 and 0A2J.
\url{https://stacks.math.columbia.edu/tag/03R1}.
\url{https://stacks.math.columbia.edu/tag/0A2J}.
\bibitem{MilneAbelianVarieties2008}
J.~S.~Milne, \emph{Abelian Varieties}, version 2.00, 2008,
I, \S3, Theorem~3.2 and Proposition~3.10.
\url{https://www.jmilne.org/math/CourseNotes/AV.pdf}.

\bibitem{BuchweitzPavlovMoore}
R.-O.~Buchweitz and A.~Pavlov,
\emph{Moore Matrices and Ulrich Bundles on an Elliptic Curve},
arXiv:1511.05502, Theorem~A and Remark~1.4(b).
\url{https://arxiv.org/abs/1511.05502}.

\bibitem{KollarCubicUnirational}
J.~Koll\'ar, \emph{Unirationality of cubic hypersurfaces},
J. Inst. Math. Jussieu \textbf{1} (2002), no.~3, 467--476, Theorem~2.
\url{https://doi.org/10.1017/S1474748002000117}.
\url{https://arxiv.org/abs/math/0005146}.
\end{thebibliography}
\end{document}
